% Quantitative Finance submission, set in the journal's rQUF2e class.

\RequirePackage[2018-12-01]{latexrelease}
\documentclass{rQUF2e}

% Additional packages (not loaded by rQUF2e)
\usepackage{cmap}
\usepackage{threeparttable}
\usepackage{booktabs}
\usepackage{bm}
\usepackage{url}

% Theorem environments (as per rQUFguide)
\theoremstyle{plain}
\newtheorem{theorem}{Theorem}[section]
\newtheorem{corollary}[theorem]{Corollary}
\newtheorem{lemma}[theorem]{Lemma}
\newtheorem{proposition}[theorem]{Proposition}

\theoremstyle{definition}
\newtheorem{definition}{Definition}

\theoremstyle{remark}
\newtheorem{remark}{Remark}

% Auto-generated numbers for prose references (regenerate with code/output/generate_table_bodies.py)
% Auto-generated by code/output/generate_table_bodies.py
% Do not edit — regenerate with: make tables
\newcommand{\persistMinDays}{250}
\newcommand{\dirNullDraws}{10{,}000}
\newcommand{\dirNullCap}{0.053}
\newcommand{\dirExcessCap}{1.54}
\newcommand{\dirExcessFor}{0.69}
\newcommand{\dirNullRatio}{1.2}
\newcommand{\dirExcessRatio}{2.2}
\newcommand{\dirExcessCapPre}{2.01}
\newcommand{\dirExcessForPre}{0.88}
\newcommand{\dirExcessRatioPre}{2.3}
\newcommand{\dirNullP}{0.0002}
\newcommand{\dirGaussianCap}{0.033}
\newcommand{\dirNullTwSitcP}{0.288}
\newcommand{\dirNullTwSitcExcessShare}{0.64}
\newcommand{\dirNullTwForExcessShare}{0.53}
\newcommand{\dirNullTwSitcPShare}{0.0002}
\newcommand{\rpHRMCtone}{8.71}
\newcommand{\rpHRTVtone}{-5.93}
\newcommand{\mcollCorr}{0.552}
\newcommand{\mcollVif}{1.477}
\newcommand{\oosPPMChr}{10.45}
\newcommand{\oosPPTVhr}{-7.86}
\newcommand{\sampPPDays}{1226}
\newcommand{\ryMinHRMCt}{2.25}
\newcommand{\ryMaxHRMCt}{5.37}
\newcommand{\ryMinHRMCYear}{2025}
\newcommand{\ryMaxHRMCYear}{2021}
\newcommand{\ryYearsCovered}{2020--2025}
\newcommand{\ryNYears}{6}
\newcommand{\lsMCbps}{12.98}
\newcommand{\lsMCt}{8.70}
\newcommand{\lsTVbps}{7.38}
\newcommand{\fnAlpha}{13.50}
\newcommand{\fnAlphaT}{9.35}
\newcommand{\fnSMB}{-0.224}
\newcommand{\fnSMBt}{-5.04}
\newcommand{\ihAggregateMCt}{9.33}
\newcommand{\ihAggregateTVt}{2.37}
\newcommand{\ihAggregateHRMCt}{9.74}
\newcommand{\ihAggregateHRTVt}{-4.25}
\newcommand{\winMCt}{9.11}
\newcommand{\winTVt}{-6.96}
\newcommand{\nwFiveMCt}{8.66}
\newcommand{\nwTenMCt}{8.72}
\newcommand{\nwFiveTVt}{-5.52}
\newcommand{\nwTenTVt}{-5.56}
\newcommand{\ffPPMChrt}{9.25}
\newcommand{\ffPPTVhrt}{19.40}
\newcommand{\fryMinHRTVt}{4.32}
\newcommand{\fryMaxHRTVt}{10.53}
\newcommand{\fryMinHRTVYear}{2024}
\newcommand{\fryMaxHRTVYear}{2021}
\newcommand{\fryYearsCovered}{2020--2025}
\newcommand{\ffBaseMChrt}{-1.92}
\newcommand{\ffBaseTVhrt}{18.66}
\newcommand{\ffWinMChrt}{0.61}
\newcommand{\ffWinTVhrt}{17.14}
\newcommand{\fNwTenTVt}{17.66}
\newcommand{\fLsTVbps}{21.58}
\newcommand{\fLsTVt}{18.55}
\newcommand{\fFnAlpha}{21.46}
\newcommand{\fFnAlphaT}{17.26}
\newcommand{\liqCapSizeQOneT}{0.60}
\newcommand{\liqCapSizeQTwoT}{0.53}
\newcommand{\liqForSizeQFiveT}{0.64}
\newcommand{\liqForSizeQFiveMisT}{1.41}
\newcommand{\liqCapSizeQOneDays}{102}
\newcommand{\liqCommonSizeQFivePct}{54.2}
\newcommand{\liqForSizeQFivePct}{20.1}
\newcommand{\descTurnMean}{2.2}
\newcommand{\descTurnStd}{9.6}
\newcommand{\revCommonStockDays}{1{,}103{,}461}
\newcommand{\revCommonRev}{5.00}
\newcommand{\revCommonT}{6.92}
\newcommand{\revCommonCapEdge}{3.63}
\newcommand{\revCommonCapEdgeT}{8.15}
\newcommand{\revCommonForEdge}{-1.37}
\newcommand{\revCommonForEdgeT}{-2.27}
\newcommand{\revUnionRev}{8.36}
\newcommand{\revUnionCapEdge}{3.68}
\newcommand{\revUnionForEdge}{-4.68}
\newcommand{\revZeroRev}{6.91}
\newcommand{\revZeroT}{14.20}
\newcommand{\revZeroCapEdge}{2.26}
\newcommand{\revZeroCapEdgeT}{8.60}
\newcommand{\revZeroForEdge}{-4.66}
\newcommand{\revZeroForEdgeT}{-10.94}
\newcommand{\revCommonClusterP}{2.3 \times 10^{-14}}
\newcommand{\revUnionRows}{4{,}403{,}344}
\newcommand{\revUnionClusterT}{13.69}
\newcommand{\revPreCommonStockDays}{794{,}013}
\newcommand{\revPreCommonRev}{5.10}
\newcommand{\revPreCommonT}{7.38}
\newcommand{\revPreCommonCapEdge}{4.40}
\newcommand{\revPreCommonCapEdgeT}{10.39}
\newcommand{\revPreCommonForEdge}{-0.69}
\newcommand{\revPreCommonForEdgeT}{-1.23}
\newcommand{\revPreUnionRev}{5.19}
\newcommand{\revPreUnionT}{7.78}
\newcommand{\revPreUnionForEdge}{-0.71}
\newcommand{\revPreUnionForEdgeT}{-1.40}
\newcommand{\revSkipCommonRev}{-0.18}
\newcommand{\revSkipCommonT}{-0.28}
\newcommand{\revSkipCommonForEdgeT}{2.00}
\newcommand{\revSkipCommonCapEdge}{0.91}
\newcommand{\revSkipCommonForEdge}{1.09}
\newcommand{\revSkipCommonCapEdgeT}{2.39}
\newcommand{\revSkipUnionRev}{0.73}
\newcommand{\revSkipUnionT}{1.30}
\newcommand{\revSkipUnionForEdgeT}{0.45}
\newcommand{\revZeroRows}{6{,}641{,}588}
\newcommand{\revTrustStockDays}{1{,}520{,}445}
\newcommand{\revTrustRev}{5.64}
\newcommand{\revTrustT}{8.32}
\newcommand{\revTrustCapEdge}{3.24}
\newcommand{\revTrustCapEdgeT}{6.43}
\newcommand{\revTrustForEdge}{-2.40}
\newcommand{\revTrustForEdgeT}{-4.27}
\newcommand{\advWindow}{20}
\newcommand{\advWindowMin}{5}
\newcommand{\advWindowMax}{60}
\newcommand{\advSampleShare}{95.2}
\newcommand{\advBaseRev}{4.92}
\newcommand{\advBaseT}{7.03}
\newcommand{\advRev}{3.25}
\newcommand{\advT}{5.00}
\newcommand{\advClusterT}{5.15}
\newcommand{\advJointRev}{9.01}
\newcommand{\advJointT}{5.21}
\newcommand{\advCapEdge}{1.82}
\newcommand{\advCapEdgeT}{4.92}
\newcommand{\advCapEdgeBase}{3.11}
\newcommand{\advForEdge}{-1.43}
\newcommand{\advForEdgeT}{-2.66}
\newcommand{\advForEdgeBase}{-1.81}
\newcommand{\advGap}{1.67}
\newcommand{\advSweepRevLo}{2.98}
\newcommand{\advSweepRevHi}{3.74}
\newcommand{\advSweepTLo}{4.65}
\newcommand{\advUnionRev}{4.18}
\newcommand{\advUnionT}{7.97}
\newcommand{\cemIntensityCap}{1.2}
\newcommand{\cemIntensityFor}{6.9}
\newcommand{\cemBalanceBar}{0.1}
\newcommand{\cemGapFull}{0.91}
\newcommand{\cemGapMatched}{0.10}
\newcommand{\cemGapDecile}{0.03}
\newcommand{\cemMagGapFull}{0.88}
\newcommand{\cemMagGapMatched}{0.14}
\newcommand{\cemMagGapDecile}{0.06}
\newcommand{\cemEssCapPct}{74}
\newcommand{\cemEssForPct}{81}
\newcommand{\cemRev}{4.36}
\newcommand{\cemT}{4.73}
\newcommand{\cemOverlapRev}{5.01}
\newcommand{\cemOverlapT}{6.86}
\newcommand{\cemCapEdge}{3.72}
\newcommand{\cemCapEdgeT}{7.24}
\newcommand{\cemForEdge}{-0.63}
\newcommand{\cemForEdgeT}{-0.79}
\newcommand{\propBootDraws}{2{,}000}
\newcommand{\contCapMC}{+0.322}
\newcommand{\contCapTV}{+0.396}
\newcommand{\contCapMCT}{+119.09}
\newcommand{\contCapTVT}{+179.98}
\newcommand{\propCapCross}{+1.59}
\newcommand{\propCapCrossLo}{+1.22}
\newcommand{\propCapCrossHi}{+1.94}
\newcommand{\propCapCrossZ}{+8.54}
\newcommand{\propCapQMC}{+0.226}
\newcommand{\propCapQTV}{+0.396}
\newcommand{\propCapQMCT}{+144.02}
\newcommand{\propCapQTVT}{+179.98}
\newcommand{\propCapCrossCap}{+0.76}
\newcommand{\propCapCrossCapZ}{+7.22}
\newcommand{\contForMC}{-0.066}
\newcommand{\contForTV}{+0.172}
\newcommand{\contForMCT}{-23.46}
\newcommand{\contForTVT}{+63.26}
\newcommand{\propForCross}{+0.36}
\newcommand{\propForCrossLo}{+0.20}
\newcommand{\propForCrossHi}{+0.53}
\newcommand{\propForCrossZ}{+4.25}
\newcommand{\propForQMC}{+0.010}
\newcommand{\propForQTV}{+0.172}
\newcommand{\propForQMCT}{+6.77}
\newcommand{\propForQTVT}{+63.26}
\newcommand{\propForCrossCap}{+0.57}
\newcommand{\propForCrossCapZ}{+14.05}
\newcommand{\fexGridLo}{-3.0}
\newcommand{\fexGridHi}{+4.0}
\newcommand{\fexGridStep}{0.1}
\newcommand{\fexBootDraws}{2{,}000}
\newcommand{\fexCapAssignedPct}{64.2}
\newcommand{\fexCapOtherPct}{8.3}
\newcommand{\fexRankCapAssignedPct}{71.7}
\newcommand{\fexRankCapOtherPct}{33.6}
\newcommand{\fexForAssignedPct}{99.1}
\newcommand{\fexForOtherPct}{32.4}
\newcommand{\fexRankForAssignedPct}{95.1}
\newcommand{\fexRankForOtherPct}{80.2}
\newcommand{\fexGap}{+2.20}
\newcommand{\fexGapLo}{+1.80}
\newcommand{\fexGapHi}{+4.10}
\newcommand{\fexCommonForAssignedPct}{90.6}
\newcommand{\fexCommonForOtherPct}{42.9}
\newcommand{\fexForRhoMax}{0.0220}
\newcommand{\fexCommonForRhoMax}{0.0110}
\newcommand{\fexCommonGap}{+2.50}
\newcommand{\fexCommonGapLo}{-2.10}
\newcommand{\fexCommonGapHi}{+4.40}
\newcommand{\sizSmallQuintiles}{3}
\newcommand{\sizLargeQuintiles}{2}
\newcommand{\sizSmallStockDays}{195{,}401}
\newcommand{\sizSmallPerDay}{114}
\newcommand{\sizLargeStockDays}{908{,}060}
\newcommand{\sizLargePerDay}{623}
\newcommand{\sizSmallSharePct}{17.7}
\newcommand{\sizSmallLvCap}{+5.14}
\newcommand{\sizSmallLvCapT}{+4.69}
\newcommand{\sizSmallLvFor}{-5.87}
\newcommand{\sizSmallLvForT}{-3.22}
\newcommand{\sizSmallLvRev}{+11.00}
\newcommand{\sizSmallLvRevT}{+5.39}
\newcommand{\sizSmallRkCap}{+3.35}
\newcommand{\sizSmallRkCapT}{+4.74}
\newcommand{\sizSmallRkFor}{-1.51}
\newcommand{\sizSmallRkForT}{-1.38}
\newcommand{\sizSmallRkRev}{+4.86}
\newcommand{\sizSmallRkRevT}{+4.03}
\newcommand{\sizSmallCemCap}{+5.15}
\newcommand{\sizSmallCemCapT}{+3.84}
\newcommand{\sizSmallCemFor}{-2.19}
\newcommand{\sizSmallCemForT}{-0.99}
\newcommand{\sizSmallCemRev}{+7.34}
\newcommand{\sizSmallCemRevT}{+3.03}
\newcommand{\sizLargeLvCap}{+3.56}
\newcommand{\sizLargeLvCapT}{+7.57}
\newcommand{\sizLargeLvFor}{-0.31}
\newcommand{\sizLargeLvForT}{-0.53}
\newcommand{\sizLargeLvRev}{+3.87}
\newcommand{\sizLargeLvRevT}{+5.34}
\newcommand{\sizLargeRkCap}{+2.40}
\newcommand{\sizLargeRkCapT}{+9.40}
\newcommand{\sizLargeRkFor}{+0.46}
\newcommand{\sizLargeRkForT}{+1.32}
\newcommand{\sizLargeRkRev}{+1.94}
\newcommand{\sizLargeRkRevT}{+4.70}
\newcommand{\sizLargeCemCap}{+3.61}
\newcommand{\sizLargeCemCapT}{+6.71}
\newcommand{\sizLargeCemFor}{+0.37}
\newcommand{\sizLargeCemForT}{+0.51}
\newcommand{\sizLargeCemRev}{+3.23}
\newcommand{\sizLargeCemRevT}{+3.64}
\newcommand{\sizResampleDraws}{40}
\newcommand{\sizResampleMean}{-1.35}
\newcommand{\sizResampleSd}{1.21}
\newcommand{\sizResampleZ}{3.74}
\newcommand{\sizResampleSigPct}{12.5}
\newcommand{\sizSmallGapLo}{-1.00}
\newcommand{\sizLargeGapLo}{-2.30}
\newcommand{\rankCapMCt}{9.45}
\newcommand{\rankCapTVt}{5.77}
\newcommand{\rankCapHRMCt}{10.86}
\newcommand{\rankCapHRTVt}{-9.11}
\newcommand{\rankForMCt}{13.68}
\newcommand{\rankForTVt}{18.90}
\newcommand{\rankForHRMCt}{2.03}
\newcommand{\rankForHRTVt}{11.18}
\newcommand{\rankCommonRev}{2.42}
\newcommand{\rankCommonT}{5.97}
\newcommand{\rankCommonCapEdge}{2.44}
\newcommand{\rankCommonCapEdgeT}{10.16}
\newcommand{\rankCommonForEdge}{0.02}
\newcommand{\rankCommonForEdgeT}{0.05}
\newcommand{\rankUnionRev}{3.59}
\newcommand{\rankUnionT}{10.20}
\newcommand{\rankUnionCapEdge}{2.45}
\newcommand{\rankUnionCapEdgeT}{10.17}
\newcommand{\rankUnionForEdge}{-1.14}
\newcommand{\rankUnionForEdgeT}{-4.14}
\newcommand{\rankCorrCapZ}{0.56}
\newcommand{\rankCorrCapRank}{0.88}
\newcommand{\rankCorrForZ}{0.47}
\newcommand{\rankCorrForRank}{0.78}
\newcommand{\doseClasses}{9}
\newcommand{\doseIndAdv}{-0.37}
\newcommand{\doseIndMC}{-7.37}
\newcommand{\doseIndTV}{-7.00}
\newcommand{\doseIndAdvT}{-0.65}
\newcommand{\doseBankT}{0.46}
\newcommand{\dosePenAdv}{3.63}
\newcommand{\dosePenAdvT}{7.57}
\newcommand{\dosePenObsApprox}{0.91~million}
\newcommand{\doseInsAdv}{2.71}
\newcommand{\doseInsAdvT}{5.91}
\newcommand{\doseInsObsApprox}{0.84~million}
\newcommand{\vwapCapMatchedClose}{8.71}
\newcommand{\vwapCapMatchedVwap}{8.62}
\newcommand{\vwapCapMismatchClose}{-5.93}
\newcommand{\vwapCapMismatchVwap}{-5.90}
\newcommand{\vwapForMatchedClose}{18.66}
\newcommand{\vwapForMatchedVwap}{18.73}
\newcommand{\vwapForMismatchClose}{-1.92}
\newcommand{\vwapForMismatchVwap}{-2.08}
\newcommand{\sampTotalDays}{1{,}473}
\newcommand{\sampTotalDaysMinusOne}{1{,}472}
\newcommand{\sampFactorDays}{1{,}472}
\newcommand{\panelTotalObsApprox}{3.32~million}
\newcommand{\sampForeignObs}{3{,}298{,}239}
\newcommand{\sampInstPanelPct}{33.4}
\newcommand{\sampForeignPanelPct}{99.2}
\newcommand{\foreignOtherGrossPct}{1.0}
\newcommand{\foreignOtherFlipPct}{1.2}
\newcommand{\foreignZeroAbsentPct}{89.0}
\newcommand{\capZeroAbsentPct}{99.4}
\newcommand{\twnPanelRows}{2{,}571{,}372}
\newcommand{\twnPanelStocks}{1{,}974}
\newcommand{\twnYears}{2020--2025}
\newcommand{\twnCorpActionDrift}{0.30}
\newcommand{\twnCommonCapEdge}{7.85}
\newcommand{\twnCommonCapEdgeT}{14.52}
\newcommand{\twnCommonForEdge}{3.95}
\newcommand{\twnCommonForEdgeT}{6.41}
\newcommand{\twnCommonRev}{3.90}
\newcommand{\twnCommonT}{4.64}
\newcommand{\twnUnionCapEdge}{7.85}
\newcommand{\twnUnionCapEdgeT}{14.47}
\newcommand{\twnUnionForEdge}{1.91}
\newcommand{\twnUnionForEdgeT}{5.51}
\newcommand{\twnUnionRev}{5.94}
\newcommand{\twnUnionT}{9.02}
\newcommand{\twnSurvivorStocks}{6}
\newcommand{\twnSurvivorSessions}{263}
\newcommand{\retForeignKospiAllObs}{3{,}298{,}239}
\newcommand{\retForeignKospiInObs}{294{,}279}
\newcommand{\retForeignKospiOutObs}{3{,}003{,}960}
\newcommand{\retForeignKospiInDailyAvg}{200}
\newcommand{\retForeignKospiOutDailyAvg}{2039}
\newcommand{\minStocksPerDay}{30}
\newcommand{\sampleYears}{2020--2025}
\newcommand{\preYears}{2015--2019}
\newcommand{\nwLagsFive}{5}
\newcommand{\nwLagsTen}{10}
\newcommand{\winsorLowerPct}{1}
\newcommand{\winsorUpperPct}{99}
\newcommand{\threshGridKRW}{0, 100, 250, 500, 1000, and 2000}
\newcommand{\threshTopRemovedPct}{81}
\newcommand{\threshTopTickers}{318}
\newcommand{\threshTopMCt}{1.91}
\newcommand{\threshTopTVt}{+0.21}
\newcommand{\threshHeldRemovedPct}{26}
\newcommand{\threshHeldMCt}{8.08}
\newcommand{\threshMCLastPct}{68}
\newcommand{\threshMCLastT}{3.84}
\newcommand{\threshTopCoefRatio}{0.43}
\newcommand{\threshTopSeRatio}{1.97}
\newcommand{\threshLsBpsMCMin}{12.3}
\newcommand{\threshLsBpsMCMax}{13.4}
\newcommand{\threshLsBpsTVMin}{7.4}
\newcommand{\threshLsBpsTVMax}{9.3}
\newcommand{\threshIcMCmaxT}{0.89}
\newcommand{\rySignifNegYears}{2021, 2022, 2023, and 2024}
\newcommand{\ryInsignifYears}{2020 and 2025}

% The manuscript's own word count (regenerated by make pdf, from code/output/word_count.py).
% TeXcount skips the input, so the count never depends on the file recording it.
%TC:ignore
% Auto-generated by code/output/word_count.py -- regenerate with: make pdf
\newcommand{\manuscriptWordCount}{12{,}816}

%TC:endignore

% rQUF2e records the last page number before the floats still waiting at the
% end of the document are set, so the first page's page range would stop one
% page early. Setting those floats first, and then stepping back to the last
% page set, makes the recorded number the PDF's last page.
\AtEndDocument{\clearpage\addtocounter{page}{-1}}

\begin{document}

\title{Signal extraction from order flow: a matched-filter perspective on normalization}

\markboth{Quantitative Finance}{Signal extraction from order flow}

\maketitle

\begin{abstract}
No criterion chooses between market capitalization and traded value as the denominator of order flow. A matched filter, an estimator of a signal of known shape, divides an investor class's flow by the characteristic that scales the flow's return-predicting signal. That characteristic is, by assumption, the one to which the class sizes its trading: market capitalization for pension funds and insurers, one class, and traded value for foreign investors. A return-free measure of how much of each class's flow accumulates across days, as positions do, corroborates those sizing rules. A class's advantage of market-capitalization normalization is its slope of next-day returns on its flow over market capitalization minus its slope on its flow over traded value. In Korea (\sampleYears{}), the advantage is significantly positive for pension and insurance flow and significantly negative for foreign flow on stock-days that both classes trade. The contrast between the two advantages is significant. The contrast and the advantage for pension and insurance flow, though not the advantage for foreign flow, remain significant on a \preYears{} pre-period and in the top \sizLargeQuintiles{} size quintiles. Market-capitalization normalization, common in studies of informed institutional trading, thus suits pension and insurance flow better than it suits foreign flow.
\end{abstract}

\begin{keywords}
Signal extraction; Order flow; Matched filter; Market microstructure; Normalization; Investor heterogeneity
\end{keywords}

\begin{classcode}G12, G14, C58\end{classcode}

{\keywordfont\noindent{\itshape Word count\/}: \manuscriptWordCount{}, excluding references and table entries.\par}


\section{Introduction}


Order flow is the net buying of a stock by a group of investors. Studies normalize it by dividing net buying by a measure of firm size or of trading activity, whether they estimate the probability of informed trading \citep{Easley1996}, construct liquidity factors \citep{AMIHUD200231, PastorStambaugh2003}, or design execution algorithms \citep{Almgren2003}. Studies of informed institutional trading tend to divide by market capitalization or shares outstanding \citep{CAMPBELL200966, Beber2011, LEWELLEN201162}, whereas studies of liquidity, market impact, and flow toxicity (the adverse selection that informed flow imposes on liquidity providers) divide by traded value, the currency value of the day's trades \citep{AMIHUD200231, Easley2012}. Studies of order imbalance, buyer-initiated less seller-initiated trading, divide it by the number of trades or by traded value to eliminate the effect of total trading activity \citep{ChordiaSubrahmanyam2004}. Accounting research has measured how dividing a regression's variables by a size variable biases the coefficients and how a dependence on size remains after the division \citep{BarthKallapur1996, EastonSommers2003}. Each of the studies cited so far treats firm size or trading activity as a nuisance to be removed. None asks which firm characteristic the return-predicting part of an investor class's flow is proportional to, so none supplies a testable criterion for choosing between two normalizations. None says in which direction the wrong choice distorts the normalized flow.

How the flow of an investor class relates to prices is well documented, most directly in Korea, the market that this paper studies. \citet{Choe1999} document positive-feedback trading and herding by foreign investors in Korea around the 1997 crisis. \citet{Choe2005} show that foreign investors in the same market execute at systematically less favorable intraday prices than domestic investors. \citet{Richards2005} finds that daily net foreign flows into Korea and other Asian emerging equity markets have larger price impacts than earlier estimates implied. \citet{Choe1999}, \citet{Choe2005}, and \citet{Richards2005} ask whether investor-class flow moves or predicts prices and take the normalization of the flow as given. The question that precedes any such measurement is which firm characteristic each investor class's flow should be divided by.

We treat order flow normalization as a signal-extraction problem and draw on matched-filter theory from communications engineering \citep{Turin1960}. A matched filter estimates a signal of known shape by correlating the received data with that shape and dividing the correlation by the shape's squared norm. In order flow, the signal is the information in a class's flow that predicts returns, and the rest of the flow is noise. The signal's known shape is its carrier, the firm characteristic in proportion to which the flow carries the signal: market capitalization for a class that sizes its positions in proportion to market capitalization, and traded value for a class that sizes its trades in proportion to the day's traded value. Each stock gives one observation of the flow per day, so the correlation is the flow times the carrier, and dividing the correlation by the carrier's square leaves the flow divided by the carrier. Traded value equals market capitalization times turnover, so division by the other characteristic, the mismatched filter, multiplies the signal by the stock's turnover or by its inverse and makes the normalized flow vary with turnover even where the signal does not. The paper models the noise as proportional to the other characteristic, so the matched filter makes the noise vary with turnover instead. Which filter predicts returns better therefore depends on how strong the signal is relative to the noise, because the distortion of the signal matters most where the signal dominates the flow and the variation of the noise matters most where the noise dominates. Under the model of Section~\ref{sec:theory}, the matched filter predicts better above a threshold of that signal-to-noise ratio, but the threshold's derivation assumes that the signal and the noise are independent of turnover, an assumption that order flow cannot satisfy, because net buying cannot exceed traded value (Section~\ref{subsec:threshold_estimate}). The paper therefore tests on returns whether each class's matched filter predicts better.

For two investor classes whose signals have different carriers, the matched filter of one is the mismatched filter of the other. If each class's matched filter predicts better, two such classes rank the two normalizations in opposite orders, a reversal. Capacity-bound institutions, the pension funds and insurers whose mandates and capacity constraints tie position size to market capitalization, should predict returns better when their flow is divided by market capitalization than when it is divided by traded value, whereas foreign flow should reverse that ordering. The paper measures both orderings on the stock-days that both classes trade, the same firms, days, and prices, so a confound that favors one normalization equally in both classes cannot make one class favor that normalization more than the other.

Each class's filter is assigned from the class's scaling role, the sizing rule of its documented trading: positions sized to market capitalization or trades sized to traded value. Pension funds and insurers accumulate positions sized to the firm's market capitalization, whereas foreign investors execute through trades sized to the day's traded value \citep{Grinold1999, LEWELLEN201162, Almgren2003}. A statistic that uses no returns corroborates the assignment, so the reversal tests the assignment rather than fitting it to the returns. That statistic is the directionality ratio, the absolute value of a stock's mean daily flow over the mean of its absolute daily flow. The ratio is high for a class whose purchases and sales accumulate in one direction and low for a class whose purchases and sales offset over time. The paper takes a higher ratio to mark a class that sizes positions to market capitalization, and a lower ratio to mark a class that sizes trades to the day's traded value. The ratio is higher for capacity-bound institutions than for foreign investors on the analysis window (\sampleYears{}) and on the pre-period (\preYears{}), which precedes it. A class whose daily flows had random signs would still have a positive ratio, which rises as the number of days on which the class trades falls. Capacity-bound institutions trade on fewer days than foreign investors, so the paper also divides each class's ratio by its value under random signs of the same daily flows, and the capacity-bound ratio remains the higher.

A class's advantage of market-capitalization normalization is its slope of next-day returns on its market-capitalization-normalized flow minus the slope on its traded-value-normalized flow, each flow standardized within the day, so that the two slopes rank the two filters as their correlations with returns do. The contrast, the paper's test of the reversal, is the advantage for capacity-bound flow minus the advantage for foreign flow. The contrast has $t = \revCommonT$ on the \revCommonStockDays{} stock-days that both classes trade and that have a next-day return. Each class's advantage also has the predicted sign, positive for capacity-bound flow and negative for foreign flow, and is significant on its own on those stock-days (Section~\ref{subsec:reversal_test}). The contrast and the advantage for capacity-bound flow remain significant at the 1\% level in four checks of whether the sample period, firm size, a few extreme flows, or a difference in how heavily the two classes trade produces them: on the pre-period, in the top \sizLargeQuintiles{} size quintiles, under a within-day rank transform of the normalized flows, and when the two classes are matched on the size of their net buying relative to traded value (Section~\ref{subsec:contrast_halves}). On the stock-days that both classes trade, the advantage for foreign flow is significant at the 5\% level in none of the four checks, so the four checks show that the two classes differ and do not show that foreign flow reverses the ordering. The contrast is limited to the next day's return, because the contrast on the return over the second day after the flow is indistinguishable from zero. In a second market, Taiwan, where investment trusts tie position size to market capitalization as pension funds and insurers do, the contrast between investment-trust and foreign flow is also positive and significant on the stock-days that both classes trade, although the advantage for foreign flow there is positive (Section~\ref{subsec:second_market}). Market-capitalization normalization, which studies of informed institutional trading tend to use \citep{CAMPBELL200966, Beber2011, LEWELLEN201162}, therefore suits capacity-bound flow better than it suits foreign flow.

The paper contributes a criterion that assigns the denominator of an investor class's order flow from the class's scaling role and checks the assignment with a statistic that uses no returns. The paper also tests the criterion in two markets, Korea and Taiwan, whose exchanges disclose each investor class's net buying daily and so separate two classes that the criterion assigns different denominators.


%%%%%%%%%%%%%%%%%%%%%%%%%%%%%%%%%%%%%%%%%%%%%%%%%%%%%%%%%%%%%%%%%%%%%%%%%%%%%%%
\section{A matched-filter model of order flow normalization}
\label{sec:theory}
%%%%%%%%%%%%%%%%%%%%%%%%%%%%%%%%%%%%%%%%%%%%%%%%%%%%%%%%%%%%%%%%%%%%%%%%%%%%%%%

\subsection{Setup and primitives}

We formalize the normalization problem as a signal-extraction problem, applying the matched-filter principle from communications engineering \citep{Turin1960}. A receiver that detects a signal of known shape in noise estimates the signal's amplitude with a matched filter, which correlates the received data with a template of that shape and divides the correlation by the template's squared norm. In order flow, the shape is the signal's carrier, the firm characteristic to which the flow's signal component is proportional (Section~\ref{subsec:general_model}). Each stock gives one observation, so the estimate is the stock's flow divided by that characteristic. Consider a cross-section of $N$ stocks indexed by $i \in \{1,\ldots,N\}$. Each stock has market capitalization $M_i$ (known, time-invariant for simplicity) and daily traded value $V_i = \tau_i M_i$ (in currency units), where $\tau_i$ is the turnover rate.

\subsection{General order flow model}
\label{subsec:general_model}

We model an investor class's observed order flow as the sum of a signal component and a noise component, each proportional to a firm characteristic.
\begin{equation}
D_i = k \cdot s_i \cdot G_i + \nu_i \cdot H_i
\label{eq:general_dgp}
\end{equation}
where $D_i$ is observed order flow (in currency units), $s_i$ is the latent signal, the information in the flow that predicts returns through equation~\eqref{eq:returns}, $\nu_i$ is noise (independent of $s_i$, $\tau_i$, $M_i$), and $k > 0$ is a scaling constant. The signal and the noise have mean zero and variances $\sigma_s^2$ and $\sigma_\nu^2$. The signal carrier $G_i$ and the noise carrier $H_i$ are the firm characteristics to which the signal component and the noise component are proportional. Future returns are
\begin{equation}
R_i = \gamma s_i + \epsilon_i
\label{eq:returns}
\end{equation}
where $\gamma > 0$ is the return sensitivity to the signal and $\epsilon_i$ is idiosyncratic return noise with mean zero and variance $\sigma_\epsilon^2$, independent of $(s_i, \nu_i, \tau_i, M_i)$. The model admits two economically distinct regimes, which differ in the characteristic that is the signal carrier.

\textbf{Capacity regime} ($G_i = M_i$, $H_i = V_i$). Capacity-bound institutions (pension funds and insurers) size positions relative to firm market capitalization under capital allocation rules and capacity constraints \citep{Grinold1999, LEWELLEN201162}. The noise is the part of the class's flow that responds to the day's liquidity and is proportional to traded value. Here the signal $s_i$ is fundamental information. Observed order flow is
\begin{equation*}
D_i = k s_i M_i + \nu_i V_i.
\end{equation*}

\textbf{Volume regime} ($G_i = V_i$, $H_i = M_i$). In this regime the class's traders execute through algorithms that target a participation rate, a share of the day's traded value, rather than a position, so their trades are proportional to the day's traded value \citep{Almgren2001, Almgren2003}. Algorithmic volume-weighted average price (VWAP) strategies are the canonical example. Their prevalence in modern markets is well documented \citep{Hendershott2011}. Here the signal $s_i$ is the deviation of the target participation rate from its cross-sectional average. The deviation can carry information, because an informed trader who routes orders through a volume-targeting algorithm sets a higher participation rate in the stocks in which conviction is highest. The noise $\nu_i$ is the class's position-sized trading, which carries no information about returns. Observed order flow is
\begin{equation*}
D_i = k s_i V_i + \nu_i M_i.
\end{equation*}

\begin{remark}[Signal--turnover independence]
\label{rem:independence}
In both regimes, $s_i$ is independent of $(\tau_i, M_i)$, so the signal is uncorrelated with firm size and turnover.
\end{remark}

\subsection{Signal extraction and the matched-filter principle}

The goal is to extract the latent signal $s_i$ from observed $D_i$ to predict $R_i$. Two normalizations are available, dividing by $G_i$ (the signal carrier) or by $H_i$ (the noise carrier). The choice determines how the normalized flow depends on the signal.

\textbf{Matched filter} (normalization by the signal carrier $G_i$):
\begin{align*}
S_i^{\text{matched}} = \frac{D_i}{G_i} = k s_i + \nu_i \frac{H_i}{G_i}
\end{align*}
The signal $s_i$ enters with the same coefficient $k$ for every firm, so, with equation~\eqref{eq:returns}, the expected return is the same function of the signal component $ks_i$ of $S_i^{\text{matched}}$ in every stock. The noise term $\nu_i H_i/G_i$ is heteroskedastic but leaves that function unchanged.

\textbf{Mismatched filter} (normalization by the noise carrier $H_i$):
\begin{align}
S_i^{\text{mismatched}} = \frac{D_i}{H_i} = k s_i \frac{G_i}{H_i} + \nu_i
\label{eq:mismatched}
\end{align}
The signal $s_i$ is multiplied by the carrier ratio $\omega_i = G_i/H_i$, which differs across firms, so the expected return is a different function of the signal component $ks_i\omega_i$ of $S_i^{\text{mismatched}}$ in each stock.

Because $V_i = \tau_i M_i$, the carrier ratio is a function of turnover in both regimes. Let $S^{\mathrm{MC}} = D/M$ denote market-capitalization normalization and $S^{\mathrm{TV}} = D/V$ traded-value normalization. In the capacity regime ($G=M$, $H=V$), $S^{\mathrm{MC}} = ks + \nu\tau$ is the matched filter and $S^{\mathrm{TV}} = ks\tau^{-1} + \nu$ the mismatched filter, with $\omega = \tau^{-1}$, and in the volume regime ($G=V$, $H=M$), $S^{\mathrm{TV}} = ks + \nu\tau^{-1}$ is matched and $S^{\mathrm{MC}} = ks\tau + \nu$ mismatched, with $\omega = \tau$. To compare the filters, we compare normalized flows by their squared correlation with returns, $\rho^2(S,R) = \mathrm{Cov}^2(S,R) / \{\mathrm{Var}(S)\,\mathrm{Var}(R)\}$. The return noise $\epsilon_i$ is independent of the flow, so $\rho(S,R) = \rho(S,s)\,\rho(s,R)$, and $\rho(s,R)$ is the same for both filters. Ranking the filters by squared correlation with returns therefore ranks them by $\rho^2(S,s)$, which rises with the output signal-to-noise ratio $\rho^2(S,s)/\{1-\rho^2(S,s)\}$. A matched filter maximizes that ratio when the noise has the same variance in every observation, a condition that the noise term $\nu_i H_i/G_i$ of this model's matched filter does not meet. Because this model's matched filter does not meet that condition, Proposition~\ref{prop:matched_filter} states the ranking in terms of the signal-to-noise ratio of the primitives, $\mathrm{SNR} \equiv k^2\sigma_s^2 / \sigma_\nu^2$, which no filter changes.

\begin{proposition}[General matched-filter dominance]
\label{prop:matched_filter}
Let $\rho(S,R)$ denote the cross-sectional correlation between a normalized flow $S$ and future returns $R$, evaluated at the population moments of the joint distribution of $(s_i, \nu_i, \tau_i, M_i)$, and let turnover be non-degenerate, $\text{Var}(\tau_i) > 0$. Let $E[\tau_i^2]$ and $E[\tau_i^{-2}]$ be finite. Under the general equation~\eqref{eq:general_dgp} with returns~\eqref{eq:returns},
\begin{equation*}
\rho(D/G, R) > \rho(D/H, R)
\end{equation*}
holds if and only if the signal-to-noise ratio exceeds a regime-dependent threshold,
\begin{equation}
\frac{k^2\sigma_s^2}{\sigma_\nu^2} > \Phi(\omega), \quad \text{where } \Phi(\omega) = \frac{(E[\omega])^2\, E[\omega^{-2}] - 1}{\text{Var}(\omega)} \label{eq:threshold}
\end{equation}
and $\omega = G/H$.
\end{proposition}

Rescaling each filter to the same covariance with returns makes the filter with the smaller variance the one with the larger correlation. After that rescaling, the mismatched filter has the larger signal variance, by the factor $E[\omega^2]/(E[\omega])^2 > 1$, which favors the matched filter, whereas the matched filter has the larger noise variance, by the factor $(E[\omega])^2 E[\omega^{-2}] > 1$, which favors the mismatched filter. The threshold $\Phi(\omega)$ weighs the two factors.

\begin{corollary}
\label{cor:regimes}
Under the threshold condition of Proposition~\ref{prop:matched_filter}, setting $(G,H) = (M,V)$ yields $\rho(S^{\mathrm{MC}},R) > \rho(S^{\mathrm{TV}},R)$ \textup{(capacity regime)}, and setting $(G,H) = (V,M)$ yields $\rho(S^{\mathrm{TV}},R) > \rho(S^{\mathrm{MC}},R)$ \textup{(volume regime)}.
\end{corollary}


Appendix~\ref{subsec:proof_matched_filter} proves Proposition~\ref{prop:matched_filter} and Corollary~\ref{cor:regimes}.

Equation~\eqref{eq:mismatched} is also a multiplicative errors-in-variables problem, in which the mismatched filter measures the latent signal multiplied by the firm-specific random factor $\omega_i$. The filtering formulation adds three results that the errors-in-variables reading, which takes the choice of regressor as given, does not give: a criterion for choosing between two candidate normalizations rather than a diagnosis of one, the closed-form condition~\eqref{eq:threshold} under which the better choice changes, and the duality of the two regimes (position-sized trading is the signal of one regime and the noise of the other).

\subsection{Economic interpretation}

In the capacity regime, the mismatched filter $S^{\mathrm{TV}} = D/V$ multiplies the signal by $\tau_i^{-1}$ (equation~\eqref{eq:mismatched} with $\omega = \tau^{-1}$), so among stocks with the same signal, a low-turnover stock shows a larger normalized flow, as if its informed position were larger, and a high-turnover stock a smaller normalized flow. In the volume regime, the mismatched filter $S^{\mathrm{MC}} = D/M$ multiplies the signal by $\tau_i$, so a high-turnover stock shows the larger normalized flow. In both regimes, the mismatched normalization turns differences in turnover across stocks into differences in the normalized flow that are unrelated to the signal. Reweighting the observations does not undo the distortion, because a regression of returns on the mismatched filter has fitted values that are an affine function of that filter under any weights, those of generalized least squares included, so the fitted values correlate with returns exactly as strongly as the mismatched filter does. Only a transformation of the regressor recovers the matched filter, as multiplying $S^{\mathrm{TV}}$ by $\tau_i$ in the capacity regime gives $S^{\mathrm{MC}}$.

\subsection{The independence premise}
\label{subsec:wiener}

\label{subsec:threshold_estimate}The threshold of Proposition~\ref{prop:matched_filter} rests on the premise that the signal and the noise are independent of turnover, which Remark~\ref{rem:independence} states for the signal and the definition of $\nu_i$ under equation~\eqref{eq:general_dgp} states for the noise. The multiplicative distortion of equation~\eqref{eq:mismatched} does not rest on the premise, because the equation is an algebraic identity under any joint law. Order flow cannot satisfy the premise. Net buying cannot exceed traded value, so $|S^{\mathrm{TV}}| \leq 1$, and because $S^{\mathrm{MC}} = \tau S^{\mathrm{TV}}$ identically, $|S^{\mathrm{MC}}| \leq \tau$. The component of $S^{\mathrm{MC}}$ that the premise makes independent of turnover must therefore shrink as turnover falls, whether that component is the signal, as in the capacity regime, or the noise, as in the volume regime. The threshold thus does not decide the ordering for order flow, so whether each class's matched filter predicts returns better is an empirical question, which Section~\ref{sec:empirical} answers by comparing the ordering that a return-free assignment predicts with the ordering that returns show.

Because Proposition~\ref{prop:matched_filter} ranks only the two candidate denominators, Section~\ref{subsec:boundary_extent} also estimates, for each class, the member of the family of normalizations $D/(M^{1-\theta}V^{\theta})$ that correlates best with next-day returns. The family contains the two candidate denominators at $\theta = 0$ and $\theta = 1$. Its two exponents sum to one because $D$, $M$, and $V$ are all denominated in currency, so only such exponents leave the normalized flow unchanged when a firm and its flow are scaled by a common factor.

%%%%%%%%%%%%%%%%%%%%%%%%%%%%%%%%%%%%%%%%%%%%%%%%%%%%%%%%%%%%%%%%%%%%%%%%%%%%%%%
\section{Core empirical results}
\label{sec:empirical}
%%%%%%%%%%%%%%%%%%%%%%%%%%%%%%%%%%%%%%%%%%%%%%%%%%%%%%%%%%%%%%%%%%%%%%%%%%%%%%%

\subsection{Data: Korean stock market}
\label{subsec:data}


The Korean panel covers the common stocks listed at any point over \sampleYears{} on the main market (KOSPI) of the Korea Exchange (KRX) and on its market for smaller and technology firms (KOSDAQ), including the stocks delisted during that period. Investor-class net buying and the adjusted close come from FnGuide DataGuide \citep{FnGuide2024}. Traded value and market capitalization come from the \texttt{marcap} dataset \citep{marcap2024}, which compiles the KRX's daily market data. Forward returns are computed from the adjusted close, a total-return series that corrects for corporate actions and reinvests cash dividends. A session in which a stock did not trade is dropped, so an $h$-day forward return spans the stock's next $h$ sessions in the panel. Appendix~\ref{sec:appendix_data} details the sources and universe, the investor classes and the sample restrictions. Table~\ref{tab:variable_def} defines the variables built from the data.

The KRX discloses net buying for a Domestic Institutional aggregate and for its eight sub-types (Appendix~\ref{sec:appendix_data}). Capacity-bound institutional flow is the summed net buying of two sub-types, insurers and pension funds, whose mandates tie position size to market capitalization. The KRX discloses the pension funds together with other public funds, so the class's flow includes the flow of the other public funds, which the paper assigns with the pension funds because the disclosure cannot separate them. The mandates of the investment-trust sub-type also tie position size to market capitalization, but its open-end funds also trade to meet their investors' subscriptions and redemptions \citep{CovalStafford2007}, whereas insurers and pension funds invest against long-term liabilities. The class therefore leaves the sub-type out. Appendix~\ref{subsec:appendix_aggregate_alignment} adds it and finds the contrast still significant at the 1\% level. Foreign flow is the net buying of registered foreigners, the foreign investors registered with Korea's financial regulator to trade securities. Foreign flow leaves out a small unregistered residual that the KRX also discloses (Appendix~\ref{sec:appendix_data}).

Each investor class's return regressions (Section~\ref{subsec:return_prediction}) are restricted to stock-days with nonzero net buying by that class, because a zero mostly records that the class did not trade the stock that day (Appendix~\ref{subsec:appendix_panel_filters}). Capacity-bound institutions have nonzero net buying on {1{,}108{,}586} stock-days across {2606} stocks over {1473} trading days, and foreign investors on \sampForeignObs{} stock-days. The last of those days has no next-day return, so the return regressions run on \sampTotalDaysMinusOne{} days.

% Table: Variable Definitions
\begin{table}
\begin{center}
\begin{minipage}{\textwidth}
\tbl{Variable definitions and data construction}
{\begin{tabular*}{\hsize}{\x lll@{}}\toprule
\textbf{Variable} & \textbf{Symbol} & \textbf{Definition} \\\colrule
Order flow & $D_{i,t}$ & Net buy value (KRW) $= \text{buy amount} - \text{sell amount}$ \\
Market cap & $M_{i,t}$ & Market capitalization (KRW) \\
Traded value & $V_{i,t}$ & Daily traded value (KRW), $\equiv \sum_k p_k q_k$ over intraday trades $k$ \\
MC-normalized flow & $S^{\mathrm{MC}}_{i,t}$ & Market-cap normalized flow $= D_{i,t} / M_{i,t}$ \\
TV-normalized flow & $S^{\mathrm{TV}}_{i,t}$ & Traded-value normalized flow $= D_{i,t} / V_{i,t}$ \\
Turnover & $\tau_{i,t}$ & Daily turnover rate $= V_{i,t} / M_{i,t}$ \\
Returns & $R_{i,t+h}$ & $h$-day forward simple return $= (P_{i,t+h} - P_{i,t}) / P_{i,t}$ \\\botrule
\end{tabular*}}
\tabnote{Notes: $D_{i,t}$ is computed separately for each investor class. $p_k$ and $q_k$ are the price and the quantity of trade $k$, and $P_{i,t}$ is the adjusted close. $M_{i,t}$ and $V_{i,t}$ are the daily market capitalization and traded value that the \texttt{marcap} dataset records from the KRX. Returns are close-to-close on total-return adjusted prices, which removes splits, free and paid rights offerings, and capital reductions, and reinvests the cash dividend at the ex-date price.}
\label{tab:variable_def}
\end{minipage}
\end{center}
\end{table}

\subsection{Regime assignment and the directionality ratio}
\label{subsec:regime_assignment}


Each investor class is assigned its matched filter from its scaling role, whether the class sizes its trading to market capitalization or to traded value, rather than from the returns against which the assignment is later tested. The paper takes a class in the capacity regime to accumulate more of its flow in directional positions, sized to the firm's market capitalization, than a class in the volume regime, whose trades, sized to the day's traded value, offset more over time as round trips, purchases and sales of the same stock. A statistic that uses no returns measures which of the two behaviors a class's flow shows. We average the class's daily net flow $D_t$ over the sample in two ways, for each stock that the class trades on at least \persistMinDays{} days, about one trading year, so that each average spans enough days for the class's purchases and sales to offset. The gross-flow magnitude $D^{\text{gross}} \equiv \overline{|D_t|}$ is the mean size of the daily net flow, gross of netting across days. The net-accumulation magnitude $D^{\text{net}} \equiv \bigl|\overline{D_t}\bigr|$ is the size of the mean daily net flow after netting across days. The ratio of net accumulation to gross flow is the directionality ratio, which is zero for a class whose daily flows cancel over the sample and one for a class whose daily flows all have one sign. The ratio measures how much of a class's flow accumulates in one direction, and equation~\eqref{eq:general_dgp} says nothing about accumulation (Section~\ref{subsec:future_work}), so the step from an accumulating class to a signal proportional to market capitalization, and from a class of offsetting trades to a signal proportional to traded value, is a maintained hypothesis rather than a derivation.

The directionality ratio is higher for capacity-bound institutions than for foreign flow (Table~\ref{tab:directionality_kr}), the ordering that the maintained hypothesis predicts. Both ratios are far below one, so round trips account for most of the daily flow of both classes. The gross flow, which the round trips dominate, has the larger elasticity on market capitalization for capacity-bound institutions and on traded value for foreign flow in a cross-stock regression on both characteristics (Table~\ref{tab:directionality_kr}). Each class's round trips therefore scale with the characteristic of its scaling role, the carrier that equation~\eqref{eq:general_dgp} gives the signal of the class's regime rather than the one it gives the noise (Section~\ref{subsec:future_work}). The net accumulation scales more with market capitalization for both classes. The ratio therefore separates the two classes, and the net-accumulation magnitude alone does not. The ratios are not compared with zero, because a class whose daily flows have random signs still has a positive ratio, which rises as the number of days on which the class trades falls. Capacity-bound flow is nonzero on fewer days than foreign flow, so day counts alone could make its ratio the higher.

The assignment therefore compares each class with a null of its own, the sign-flip null, which randomizes the sign of each day's flow while holding its magnitude and the stock's day count fixed. A class above its null keeps more of its gross flow as net accumulation than independent signs on the same magnitudes would, and a class below its null keeps less. A class's excess directionality, the ratio of its directionality ratio to its null, is a statistic that uses no return data (Appendix~\ref{subsec:appendix_directionality_frequency}). Capacity-bound flow has an excess directionality of \dirExcessCap{} and foreign flow of \dirExcessFor, so the two classes fall on opposite sides of the value one, as the maintained hypothesis predicts. The capacity-bound class's excess directionality exceeds the foreign class's by a factor of \dirExcessRatio{}, whereas the capacity-bound null exceeds the foreign null by a factor of only \dirNullRatio, the part of the ordering that day counts and flow magnitudes alone produce. On the pre-period, \preYears{}, which precedes the analysis window and does not overlap it, the capacity-bound class's excess directionality exceeds the foreign class's by a factor of \dirExcessRatioPre{}, at \dirExcessCapPre{} against \dirExcessForPre{} for foreign flow, and the directionality ratios themselves repeat the pattern (Panel~B of Table~\ref{tab:directionality_kr}).

% Table: Return-free regime assignment (Korea directionality decomposition)
\begin{table}
\begin{center}
\begin{minipage}{\textwidth}
\setlength{\tabcolsep}{3pt}
\tbl{Return-free regime assignment: gross flow versus net accumulation (Korea)}
{\begin{tabular*}{\hsize}{\x lcccccr@{}}\toprule
& \multicolumn{2}{c}{Gross flow ($D^{\text{gross}}$)} & \multicolumn{2}{c}{Net accumulation ($D^{\text{net}}$)} & & \\
\cmidrule(lr){2-3}\cmidrule(lr){4-5}
Investor class & $\beta_M$ & $\beta_V$ & $\beta_M$ & $\beta_V$ & Direction. & $N$ \\\colrule
\multicolumn{7}{@{}l@{}}{\textbf{Panel A: Analysis window (\sampleYears)}} \\\colrule
Capacity-bound institutions & \textbf{+1.010 (+39.2)} & +0.166 (+7.0) & \textbf{+1.154 (+27.4)} & $-0.041 (-1.0)$ & 0.082 & 1{,}146 \\
Foreign & +0.507 (+71.2) & \textbf{+0.619 (+70.1)} & \textbf{+0.794 (+33.2)} & +0.448 (+16.4) & 0.030 & 2{,}549 \\
\colrule
\multicolumn{7}{@{}l@{}}{\textbf{Panel B: Pre-period (2015--2019)}} \\\colrule
Capacity-bound institutions & \textbf{+0.663 (+21.6)} & +0.350 (+10.9) & \textbf{+0.657 (+15.4)} & +0.144 (+3.7) & 0.104 & 934 \\
Foreign & +0.573 (+70.5) & \textbf{+0.641 (+68.0)} & \textbf{+1.049 (+40.6)} & +0.335 (+11.4) & 0.039 & 2{,}095 \\\botrule

\end{tabular*}}
\tabnote{Notes: Each stock's daily net flow $D_t$ is averaged into a gross-flow magnitude $D^{\text{gross}} \equiv \overline{|D_t|}$ (the absolute value inside the average) and a net-accumulation magnitude $D^{\text{net}} \equiv \bigl|\overline{D_t}\bigr|$ (the absolute value outside it). Each is regressed across stocks on the logs of the stock's mean market capitalization $\bar M$ and mean daily traded value $\bar V$, $\log(\cdot) = a + \beta_M \log \bar M + \beta_V \log \bar V$, one observation per stock with heteroskedasticity-robust standard errors scaled by $N/(N-3)$ ($t$-statistics in parentheses). The larger $|\beta|$ in each pair is in \textbf{bold}. `Direction.' is the median directionality ratio (net accumulation over gross flow), and $N$ is the number of stocks. Flows are in KRW. Panel~B repeats the decomposition on the \preYears{} pre-period. Each row restricts to stocks with at least \persistMinDays{} nonzero-flow days.}
\label{tab:directionality_kr}
\end{minipage}
\end{center}
\end{table}

\subsection{The contrast between the two classes}
\label{subsec:return_prediction}

For both capacity-bound institutions and foreign investors, the matched filter that the assignment gives the class predicts next-day returns with a larger univariate slope than the mismatched filter (Table~\ref{tab:investor_heterogeneity}).

We estimate each slope by the Fama-MacBeth procedure \citep{fama1973risk}, which runs a cross-sectional regression on each trading day $t$ and averages the daily slopes. The daily regression of forward returns on a day-standardized normalized flow is
\begin{equation*}
R_{i,t+h} = \beta_0 + \beta_1 S_{i,t} + \epsilon_{i,t},
\end{equation*}
where $S_{i,t}$ is either $S^{\mathrm{MC}}_{i,t} = D_{i,t}/M_{i,t}$ or $S^{\mathrm{TV}}_{i,t} = D_{i,t}/V_{i,t}$, standardized within each day by subtracting the daily cross-sectional mean and dividing by the daily cross-sectional standard deviation. The estimate's standard error is the standard deviation of the daily slopes divided by the square root of the number of days. For a regressor standardized within the day, the univariate slope equals the regressor's correlation with returns, the quantity that Proposition~\ref{prop:matched_filter} ranks, times the standard deviation of the day's returns, which is the same for both filters. On each day, the univariate slopes therefore rank the filters as their correlations do, whereas the slopes with the other filter held fixed need not, so the univariate specification is the primary one. The horse race, the regression that includes $S^{\mathrm{MC}}$ and $S^{\mathrm{TV}}$ together, is reported as a check of whether the ordering holds when each filter is held fixed.

\textbf{Capacity-bound institutions.} In the horse race, the coefficient on $S^{\mathrm{MC}}$ rises and the coefficient on $S^{\mathrm{TV}}$ turns negative. No result of Section~\ref{sec:theory} predicts the sign of the mismatched filter's horse-race coefficient, and the reversal does not depend on that sign. The negative coefficient is not an imprecise estimate, because across all \sampTotalDays{} trading days of the panel the class's two filters have a correlation that averages \mcollCorr{} and a variance inflation factor that averages \mcollVif, the factor by which that correlation inflates the variance of each horse-race coefficient.

\textbf{Foreign investors.} The mismatched filter $S^{\mathrm{MC}}$ predicts next-day returns too, as the capacity-bound $S^{\mathrm{TV}}$ does, so in both classes the matched-filter principle ranks two informative filters.

\textbf{Individual investors.} The framework assigns a regime only to a class whose members share a mandate or an execution convention, and individual flow is an aggregate of households that share neither, so the principle predicts no ordering for it. Individual flow's two filters predict next-day returns with nearly equal slopes (Table~\ref{tab:investor_heterogeneity}), so the data show no ordering for individual flow either.

% Table: Investor Heterogeneity
\begin{table}
\begin{center}
\begin{minipage}{\textwidth}
\tbl{Investor heterogeneity: return prediction by investor class}
{\begin{tabular*}{\hsize}{\x lccc@{}}\toprule
& MC only & TV only & Horse race \\\colrule
\multicolumn{4}{@{}l@{}}{\textbf{Panel A: Capacity-bound institutions}} \\\colrule
$S^{\mathrm{MC}}$ & 0.0004*** & --- & 0.0006*** \\
& (7.93) & & (8.71) \\
$S^{\mathrm{TV}}$ & --- & 0.0001** & $-0.0002$*** \\
& & (2.19) & ($-5.93$) \\
$R^2$ (\%) & 0.52 & 0.19 & 0.72 \\\colrule
\multicolumn{4}{@{}l@{}}{\textbf{Panel B: Foreign investors}} \\\colrule
$S^{\mathrm{MC}}$ & 0.0003*** & --- & $-0.0001$* \\
& (4.99) & & ($-1.92$) \\
$S^{\mathrm{TV}}$ & --- & 0.0007*** & 0.0008*** \\
& & (18.05) & (18.66) \\
$R^2$ (\%) & 0.31 & 0.24 & 0.52 \\\colrule
\multicolumn{4}{@{}l@{}}{\textbf{Panel C: Individual investors}} \\\colrule
$S^{\mathrm{MC}}$ & $-0.0007$*** & --- & $-0.0005$*** \\
& ($-12.06$) & & ($-8.27$) \\
$S^{\mathrm{TV}}$ & --- & $-0.0007$*** & $-0.0005$*** \\
& & ($-19.16$) & ($-12.93$) \\
$R^2$ (\%) & 0.51 & 0.21 & 0.68 \\\botrule

\end{tabular*}}
\tabnote{Notes: Fama-MacBeth regressions of next-day returns on day-standardized normalized flows, estimated separately for each investor class and restricted to stock-days with nonzero order flow for that class. `MC only' and `TV only' are univariate, and `Horse race' includes both normalized flows. $R^2$ is the mean of the daily $R^2$, which counts each day's correlation whatever its sign, so it need not rank the filters as the slopes do. $t$-statistics in parentheses are based on Fama-MacBeth time-series standard errors. ***, **, * denote significance at the 1\%, 5\%, and 10\% levels.}
\label{tab:investor_heterogeneity}
\end{minipage}
\end{center}
\end{table}

\textbf{The reversal as a single contrast.}\label{subsec:reversal_test} The reversal is a class-by-normalization interaction. Stacking the capacity-bound and foreign samples estimates it as a contrast of two coefficients, formed on each day, with a standard error of its own. We stack the two samples into one panel with one row per (stock-day, class), standardize each normalized flow within day and class, and write the regression in its horse-race form,
\begin{equation}
\begin{split}
R_{i,t+1} = {} & a_t + c_t F + b_1 z(S^{\mathrm{MC}}_{i,t}) + b_2 z(S^{\mathrm{TV}}_{i,t}) \\
  & + b_3 F\, z(S^{\mathrm{MC}}_{i,t}) + b_4 F\, z(S^{\mathrm{TV}}_{i,t}) + \epsilon_{i,t},
\end{split}
\label{eq:reversal_interaction}
\end{equation}
where $F$ equals one on the foreign rows and zero on the capacity-bound rows, $z(\cdot)$ denotes standardization within the (day, class) cross-section, and $a_t$ and $c_t$ are the day and day-by-class intercepts. The capacity-bound coefficients are $b_1$ on $S^{\mathrm{MC}}$ and $b_2$ on $S^{\mathrm{TV}}$, and the foreign coefficients are $b_1 + b_3$ and $b_2 + b_4$. The primary univariate specification fits equation~\eqref{eq:reversal_interaction} once per filter on each day, with the other filter and its interaction dropped, so each coefficient is that filter's own cross-sectional slope.

A class's advantage of market-capitalization normalization, its advantage of $S^{\mathrm{MC}}$ over $S^{\mathrm{TV}}$, is the slope of $S^{\mathrm{MC}}$ minus the slope of $S^{\mathrm{TV}}$ for that class. The advantage for capacity-bound flow is $b_1 - b_2$, and the advantage for foreign flow is $(b_1 + b_3) - (b_2 + b_4)$. The contrast is the advantage for capacity-bound flow minus the advantage for foreign flow,
\begin{equation}
(b_1 - b_2) - \bigl[(b_1 + b_3) - (b_2 + b_4)\bigr] = b_4 - b_3 .
\label{eq:reversal_contrast}
\end{equation}
Under the assignment, each class's matched filter should predict better, so the advantage for capacity-bound flow should be positive, because $S^{\mathrm{MC}}$ is that class's matched filter. The advantage for foreign flow should be negative, because $S^{\mathrm{TV}}$ is the foreign class's matched filter. The assignment therefore predicts $b_4 > b_3$. Inference is Fama-MacBeth over the per-day contrast series, which allows the residuals of stocks on the same day to be correlated \citep{petersen2009estimating}. For the horse race, which fits both filters in one regression, we also report a pooled least-squares fit with day-by-class intercepts and with standard errors clustered by date. The pooled fit checks that the horse-race inference does not depend on averaging the daily slopes.

We estimate equation~\eqref{eq:reversal_interaction} on the \revCommonStockDays{} stock-days on which both classes trade and the stock has a next-day return, their common support, so the two classes are compared on one set of firms, days, prices, and realized returns rather than on two overlapping sets. The contrast is \revCommonRev{} basis points of next-day return per standard deviation of the normalized flow, with $t = \revCommonT$ (Table~\ref{tab:reversal_interaction}). Each class's advantage has the sign that the assignment predicts and is statistically significant on its own, at the 1\% level for capacity-bound flow and at the 5\% level for foreign flow, so on common support each class's matched filter predicts better than its mismatched filter.

The horse race and the union stacking, which stacks each class on its own nonzero-flow stock-days with a next-day return, give the contrast the same sign and significance. In the horse race, the contrast roughly doubles without a matching rise in its $t$-statistic (Table~\ref{tab:reversal_interaction}), so the horse race gives no stronger evidence for the reversal than the univariate specification. In the pooled fit, the horse-race contrast has $p = \revCommonClusterP$, so the horse race also rejects a zero contrast. The union stacking differs from common support, because restricting both classes to common support discards the foreign stock-days of the smaller stocks that no capacity-bound institution trades (Appendix~\ref{subsec:appendix_size_restricted}). The union stacking gives a sample of \revUnionRows{} rows and a larger contrast of the same sign and significance (Table~\ref{tab:reversal_interaction}), with $t = \revUnionClusterT$ for its horse-race contrast in the pooled fit.

% Table: Reversal Interaction
\begin{table}
\begin{center}
\begin{minipage}{\textwidth}
\tbl{The reversal as a single contrast}
{\begin{tabular*}{\hsize}{\x lcccc@{}}\toprule
& \multicolumn{2}{c}{Common support} & \multicolumn{2}{c}{Union} \\
& Univariate & Horse race & Univariate & Horse race \\\colrule
$z(S^{\mathrm{MC}})$ & 4.43*** & 5.81*** & 4.45*** & 5.82*** \\
& (7.88) & (8.66) & (7.93) & (8.71) \\
$z(S^{\mathrm{TV}})$ & 0.80** & $-2.45$*** & 0.77** & $-2.47$*** \\
& (2.29) & ($-5.86$) & (2.19) & ($-5.93$) \\
$F \times z(S^{\mathrm{MC}})$ & $-2.73$*** & $-5.61$*** & $-1.93$*** & $-6.85$*** \\
& ($-3.26$) & ($-5.78$) & ($-2.66$) & ($-8.34$) \\
$F \times z(S^{\mathrm{TV}})$ & 2.27*** & 5.38*** & 6.43*** & 10.12*** \\
& (4.00) & (8.11) & (12.40) & (17.43) \\
\colrule
Advantage of $S^{\mathrm{MC}}$ over $S^{\mathrm{TV}}$, capacity-bound & 3.63*** & 8.27*** & 3.68*** & 8.28*** \\
& (8.15) & (8.40) & (8.25) & (8.47) \\
Advantage of $S^{\mathrm{MC}}$ over $S^{\mathrm{TV}}$, foreign & $-1.37$** & $-2.72$** & $-4.68$*** & $-8.69$*** \\
& ($-2.27$) & ($-2.42$) & ($-10.94$) & ($-11.10$) \\
Contrast ($b_4 - b_3$) & 5.00*** & 10.99*** & 8.36*** & 16.97*** \\
& (6.92) & (7.60) & (13.95) & (13.92) \\
\colrule
Stacked rows & \multicolumn{2}{c}{2{,}206{,}922} & \multicolumn{2}{c}{4{,}403{,}344} \\
Trading days & \multicolumn{2}{c}{1{,}472} & \multicolumn{2}{c}{1{,}472} \\\botrule

\end{tabular*}}
\tabnote{Notes: Stacked estimates of equation~\eqref{eq:reversal_interaction}, one row per (stock-day, class), with both normalized flows standardized within each (day, class) cross-section. Coefficients are in basis points of next-day return per standard deviation of the normalized flow. $F$ marks the foreign rows, so $b_1$ and $b_2$ are the capacity-bound coefficients and $b_1 + b_3$, $b_2 + b_4$ the foreign coefficients. `Common support' restricts both classes to the stock-days on which both trade. `Union' stacks each class on its own nonzero-flow stock-days. `Univariate' columns fit equation~\eqref{eq:reversal_interaction} once per filter, with the other filter and its interaction dropped. `Horse race' columns fit both filters together. Every column averages per-day cross-sectional coefficients with Fama-MacBeth time-series standard errors. The two advantage rows are $b_1 - b_2$ and $(b_1 + b_3) - (b_2 + b_4)$, and the contrast row their difference, defined in equation~\eqref{eq:reversal_contrast}. $t$-statistics in parentheses. ***, **, * denote significance at the 1\%, 5\%, and 10\% levels.}
\label{tab:reversal_interaction}
\end{minipage}
\end{center}
\end{table}

\textbf{Alternative estimations.}\label{subsec:contrast_halves} The contrast and the advantage for capacity-bound flow are significant at the 1\% level under each of five alternative estimations (Table~\ref{tab:contrast_halves}). A within-day rank transform replaces each normalized flow by the standard normal quantile of its within-day rank, so a few stock-days with extreme normalized flows cannot drive a slope. A net-participation matching re-weights the two classes' stock-days within cells of date, turnover and the net participation rate $|D_{i,t}|/V_{i,t}$, which narrows the difference in how heavily the classes trade, so that the contrast compares the two classes at similar trading intensity. The panel without the nonzero-flow restriction keeps the stock-days with zero net buying, which the restriction removes far more often for the capacity-bound class than for the foreign class, so that the contrast cannot rest on which stock-days each class trades (Appendix~\ref{subsec:appendix_zero_flow}). A lagged traded-value denominator, the stock's average traded value over the \advWindow{} preceding sessions, contains none of the class's own trades of the day, whereas the day's traded value contains them, so the lagged denominator removes the overlap between the flow and the denominator of $S^{\mathrm{TV}}$ (Appendix~\ref{subsec:appendix_adv_normalisation}). Finally, the primary specification is re-estimated on the pre-period. Beyond the contrast, Appendix~\ref{sec:appendix_robustness_battery} tests each class's matched filter, mostly through its horse-race $t$-statistic, in nine robustness checks: a pre-period hold-out, yearly subsamples, a long-short portfolio, winsorization, the rank transform, autocorrelation-robust standard errors, a factor-neutral alpha, market-capitalization floors, and a formation price other than the close. Seven of the nine also test $S^{\mathrm{TV}}$, the matched filter of foreign flow, because the last two ask only whether small stocks or the closing price produce the result for the capacity-bound $S^{\mathrm{MC}}$. Each class's matched filter stays significant at the 5\% level in every check that it enters, except at a market-capitalization floor that keeps only the largest \threshTopTickers{} stocks, where the capacity-bound matched filter's horse-race $t$-statistic is not significant at the 5\% level (Appendix~\ref{subsec:appendix_j_marketcap}).

% Table: where each class's advantage and the contrast stand
\begin{table}
\begin{center}
\begin{minipage}{\textwidth}
\tbl{Each class's advantage under each alternative estimation and on the second day}
{\begin{tabular*}{\hsize}{\x lccc@{}}\toprule
& \multicolumn{2}{c}{Advantage of $S^{\mathrm{MC}}$ over $S^{\mathrm{TV}}$} & \\
\cmidrule(lr){2-3}
Specification & Capacity-bound & Foreign & Contrast \\\colrule
Primary (common support) & \revCommonCapEdge{} (\revCommonCapEdgeT) & $\revCommonForEdge$ ($\revCommonForEdgeT$) & \revCommonRev{} (\revCommonT) \\
Within-day rank transform & \rankCommonCapEdge{} (\rankCommonCapEdgeT) & \rankCommonForEdge{} (\rankCommonForEdgeT) & \rankCommonRev{} (\rankCommonT) \\
Net-participation matching & \cemCapEdge{} (\cemCapEdgeT) & $\cemForEdge$ ($\cemForEdgeT$) & \cemRev{} (\cemT) \\
No nonzero-flow restriction & \revZeroCapEdge{} (\revZeroCapEdgeT) & $\revZeroForEdge$ ($\revZeroForEdgeT$) & \revZeroRev{} (\revZeroT) \\
Lagged traded-value denominator & \advCapEdge{} (\advCapEdgeT) & $\advForEdge$ ($\advForEdgeT$) & \advRev{} (\advT) \\
Pre-period sample (\preYears{}) & \revPreCommonCapEdge{} (\revPreCommonCapEdgeT) & $\revPreCommonForEdge$ ($\revPreCommonForEdgeT$) & \revPreCommonRev{} (\revPreCommonT) \\\colrule
Second-day return & \revSkipCommonCapEdge{} (\revSkipCommonCapEdgeT) & \revSkipCommonForEdge{} (\revSkipCommonForEdgeT) & $\revSkipCommonRev$ ($\revSkipCommonT$) \\\botrule
\end{tabular*}}
\tabnote{Notes: Each row is one estimate of equation~\eqref{eq:reversal_interaction} under the univariate specification, in basis points of return per standard deviation of the normalized flow, with Fama-MacBeth $t$-statistics in parentheses. The return is the next-day return in every row except `Second-day return,' which re-estimates the primary specification on the return over the second day after the flow, $[t+1, t+2]$. The advantage of $S^{\mathrm{MC}}$ over $S^{\mathrm{TV}}$ is $b_1 - b_2$ for capacity-bound flow and $(b_1 + b_3) - (b_2 + b_4)$ for foreign flow, and the contrast is their difference, defined in equation~\eqref{eq:reversal_contrast}. Every row is estimated on common support except `No nonzero-flow restriction,' in which the two stackings of Section~\ref{subsec:reversal_test} coincide because every stock-day holds both classes. Every row is estimated on the \sampleYears{} sample except `Pre-period sample,' which re-estimates the primary specification on the disjoint \preYears{} panel. The lagged-denominator row divides the flow by the stock's average traded value over the \advWindow{} preceding sessions, on the stock-days that have at least \advWindowMax{} preceding sessions.}
\label{tab:contrast_halves}
\end{minipage}
\end{center}
\end{table}

\textbf{Horizon.}\label{subsec:skip_day} The contrast is limited to the return over the first day after the flow, $[t, t+1]$, because on the second-day return, the return over $[t+1, t+2]$, the contrast is indistinguishable from zero on both stackings (Table~\ref{tab:contrast_halves} and Appendix~\ref{subsec:appendix_skip_day}). The advantage for capacity-bound flow keeps its sign on the second day, whereas the advantage for foreign flow changes sign on common support, so on the second day $S^{\mathrm{MC}}$ predicts better than $S^{\mathrm{TV}}$ for foreign flow too.

\textbf{Other channels.}\label{sec:robustness}\label{subsec:surviving_channels} Two channels other than the flow's signal could produce the predictive power. The first is the closing price, which the market-capitalization denominator shares with the base of the next-day return, so that a closing price away from the day's other prices moves both. Forming the denominator on the day's volume-weighted average price instead of the closing price leaves the capacity-bound matched filter's predictive power in place (Appendix~\ref{subsec:vwap_formation}). The second is price pressure, the price impact of a class's trading on day $t+1$, which reaches the next-day return when the class's flow persists from day $t$. If price pressure produced all of the predictive power, each of a class's two filters would predict returns in proportion to how strongly it predicts the class's trading on day $t+1$, with one factor of proportionality, the day's price impact, for both filters. The data reject that proportionality for both classes (Appendix~\ref{subsec:appendix_flow_continuation}).

\textbf{Other denominators and classes.}\label{subsec:boundary_extent} The contrast compares one pair of classes and one pair of denominators, the $\theta = 0$ and $\theta = 1$ members of the family $D/(M^{1-\theta}V^{\theta})$ of Section~\ref{subsec:threshold_estimate}. Comparisons over more denominators and over more classes order the two classes as the contrast does. Estimated freely on each class's own nonzero-flow stock-days, the exponent that maximizes the correlation with returns is larger for foreign flow than for capacity-bound flow, by a gap whose bootstrap interval excludes zero (Appendix~\ref{subsec:free_exponent}). The comparison over classes covers \doseClasses{} investor categories: foreign flow, individual flow, and seven of the eight institutional sub-types. The eighth, the government account, has zero net buying on every stock-day of the panel. Foreign flow is the only one of the investor categories with a significantly negative advantage of $S^{\mathrm{MC}}$ over $S^{\mathrm{TV}}$, whereas each of the seven institutional sub-types with nonzero flow has a positive one, distinguishable from zero in all but one of them (Appendix~\ref{subsec:subtype_dose_response}). The evidence for the volume regime therefore comes from foreign flow alone, which is also the only category domiciled abroad, so domicile could produce the same pattern (Section~\ref{subsec:limitations}). Insurance flow and pension flow, the two categories that the capacity-bound class pools, each have a positive advantage distinguishable from zero on their own, so the advantage for capacity-bound flow does not depend on either category (Appendix~\ref{subsec:appendix_aggregate_alignment}). For the KRX domestic institutional aggregate, which contains every institutional sub-type, $S^{\mathrm{MC}}$ also predicts better than $S^{\mathrm{TV}}$.

\textbf{Taiwan.}\label{subsec:second_market} The Taiwan Stock Exchange and the Taipei Exchange disclose two classes with opposite scaling roles, foreign investors and securities investment trusts, regulated mutual funds whose capacity constraints tie position size to market capitalization. Taiwan discloses no pension or insurance class, so its investment trusts take the place of the Korean class, although their open-end funds also trade on the subscriptions and redemptions for which Section~\ref{subsec:data} leaves the Korean investment-trust sub-type out (Appendix~\ref{sec:taiwan}). In a Taiwanese panel over \twnYears{}, the investment trusts' matched filter predicts returns better than their mismatched filter, as in Korea (Table~\ref{tab:twn_return_prediction}). On the stock-days that both classes trade, the contrast is \twnCommonRev{} basis points ($t = \twnCommonT$, Table~\ref{tab:twn_contrast}), so the advantage of $S^{\mathrm{MC}}$ over $S^{\mathrm{TV}}$ is larger for the investment trusts than for foreign flow. The Taiwanese investment trusts' directionality ratio is not significantly above its own null (Table~\ref{tab:directionality_frequency}), so the comparison of a class with its own null corroborates the assignment in Korea only.

\textbf{The advantage for foreign flow.}\label{subsec:size_restricted} The advantage for foreign flow is not significant at the 5\% level under three of the five alternative estimations (Table~\ref{tab:contrast_halves}). Two of the three are the within-day rank transform and the net-participation matching on common support, whereas on the union stacking the advantage is significant under both (Appendices~\ref{subsec:appendix_j_rank} and~\ref{subsec:appendix_participation_matching}). The third is the pre-period, \preYears{}, on which the advantage has the sign that the assignment predicts but is significant at the 5\% level on neither stacking (Appendix~\ref{subsec:appendix_j_oos}), so it is significant on the analysis window only. Firm size could weaken the advantage for foreign flow on common support, because the stock-days that a capacity-bound institution trades concentrate in the largest stocks, where foreign flow predicts returns least. The bottom \sizSmallQuintiles{} of the within-day market-capitalization quintiles of the full daily cross-section hold only \sizSmallSharePct\% of common support's stock-days. The $t$-statistic of foreign $S^{\mathrm{TV}}$ declines from the smallest to the largest of those quintiles, and in the largest quintile neither foreign filter predicts returns (Appendix~\ref{sec:appendix_liquidity_stratified}). We split common support into the bottom \sizSmallQuintiles{} quintiles and the top \sizLargeQuintiles{} and re-estimate the contrast in each part under three estimations: the primary specification, the rank transform, and the net-participation matching. In the top quintiles, under all three estimations, the contrast is significant at the 1\% level and the advantage is significant for capacity-bound flow only. In the bottom quintiles, the advantage for foreign flow under the primary specification is significant at the 1\% level and larger in magnitude than on the whole of common support, whereas under the rank transform and under the matching it is not significant at the 5\% level. The rank transform shrinks that advantage by construction, because the two filters' ranks differ only through turnover and identical ranks would give an advantage of exactly zero (Appendix~\ref{subsec:appendix_j_rank}). The matching reduces the effective sample size (Appendix~\ref{subsec:appendix_participation_matching}). Neither estimation can therefore confirm or reject the advantage among smaller stocks (Appendix~\ref{subsec:appendix_size_restricted}).

\label{subsubsec:kospi200_falsification}Foreign flow does not show the volume regime's ordering everywhere, neither inside Korea's large-stock index nor in Taiwan. The KOSPI~200 is Korea's index of 200 large stocks. Foreign $S^{\mathrm{TV}}$, the matched filter that the volume regime gives foreign flow, has a larger univariate $t$-statistic outside the index than in the full foreign-flow panel, and a negative one inside the index (Table~\ref{tab:foreign_kospi200_split}), although Section~\ref{subsec:regime_assignment} assigns the volume regime to all foreign flow. Membership also tracks firm size, but the out-of-index stocks of the largest size quintile do not reproduce the negative coefficient inside the index (Appendix~\ref{subsec:appendix_kospi200}). In Taiwan, foreign flow's matched filter has neither the larger univariate slope nor the larger horse-race $t$-statistic (Table~\ref{tab:twn_return_prediction}). The Taiwanese advantage for foreign flow is positive on common support (Table~\ref{tab:twn_contrast}), so foreign flow in Taiwan does not rank the normalizations as the volume regime predicts.

%%%%%%%%%%%%%%%%%%%%%%%%%%%%%%%%%%%%%%%%%%%%%%%%%%%%%%%%%%%%%%%%%%%%%%%%%%%%%%%
\section{Discussion}
\label{sec:discussion}
%%%%%%%%%%%%%%%%%%%%%%%%%%%%%%%%%%%%%%%%%%%%%%%%%%%%%%%%%%%%%%%%%%%%%%%%%%%%%%%

\subsection{Preconditions and generalizability}

The framework needs a market that discloses the flow of a class whose members share a mandate or an execution convention. Korea discloses such classes daily, whereas in a market without that disclosure a class's scaling role is unmeasured rather than absent. Three applications lie beyond what this paper tests. First, the results concern one use of a denominator, extracting a class's signal. They leave untested a second use, the division of flow by traded value in volume-normalized measures of flow toxicity \citep{Easley2012}, which measure the adverse selection that the flow imposes on liquidity providers. Second, the assignment procedure applies to any class whose mandate gives its scaling role, including classes that this paper does not study. Comparing such a class with a class of the opposite role in the same market tests the directionality ordering that the two roles imply. Third, daily data leave intraday applications open. Within a day, a class that sizes positions to market capitalization can execute them through an algorithm that targets a participation rate, so its intraday trades are sized to volume and intraday flow could hide a signal proportional to market capitalization.

\subsection{Scope and economic significance}
\label{subsec:limitations}

Two dimensions bound the evidence, the markets in which the contrast is estimated and firm size. The contrast is estimated in two markets. A further market needs at least two disclosed classes with opposite scaling roles, because the contrast between two such classes, rather than either class's advantage alone, cancels the confounds that favor one normalization equally in both classes. US institutional flow, for example, must be inferred from the quarterly holdings that institutions report to the SEC (13F filings) or from classifying trades by size, so a test on US flow would mix the framework's predictions with classification error, whereas Korea's daily disclosure separates the flow of insurers and of pension and other public funds from the institutional aggregate. Of the two markets, only Korea supports the volume regime, through foreign flow outside the KOSPI~200 (Section~\ref{subsec:size_restricted}). Within Korea, the capacity-bound matched filter's horse-race $t$-statistic loses significance under a market-capitalization floor that keeps only the largest \threshTopTickers{} stocks (Section~\ref{subsec:return_prediction}).

Capacity-bound institutions and foreign investors differ in domicile as well as in scaling role. Foreign investors' documented herding and unfavorable intraday execution \citep{Choe1999, Choe2005} and the effect of their currency exposure could vary across stocks with turnover, the axis on which the two normalizations differ. Separating domicile from scaling role requires a market in which a domestic class occupies the volume regime.

The analysis window had short-selling bans that could change the information that a class's net selling carries. Short selling was banned in every listed stock from March 2020 to May 2021 and again from November 2023 to March 2025 \citep{FSC2020shortsale, FSC2021shortsale, FSC2023shortsale, FSC2024shortsale, FSC2025shortsale}. Between those two bans, it stayed banned in every stock outside the KOSPI~200 and the KOSDAQ~150, the large-stock indices of KOSPI and KOSDAQ \citep{FSC2021shortsale, FSC2025shortsale}. The directionality ordering, the contrast, and the advantage for capacity-bound flow each hold on the pre-period, which precedes the bans, so the bans may affect the magnitudes estimated on the analysis window but do not produce these three results.

The return tests decide which filter predicts better and are not a tradable strategy. The KRX discloses investor-class net flow after the close, so the close-to-close return measured in this paper, gross of transaction costs and market impact, overstates what a position entered after disclosure earns. The return tests also leave unmeasured the share of the predictive power that price pressure produces, because the proportionality that Section~\ref{subsec:surviving_channels} rejects holds only if pressure produces all of it.

\subsection{Model assumptions and future work}
\label{subsec:future_work}

\textbf{Signal--turnover independence.} A model whose flow components satisfy the bound $|S^{\mathrm{TV}}| \leq 1$ by construction would determine how the variances of the signal and the noise depend on turnover, which Section~\ref{subsec:threshold_estimate} leaves unknown, and so give the threshold a value.

\textbf{Single-timescale flow.} Equation~\eqref{eq:general_dgp} models order flow as a single cross-section and so merges two components that the gross-flow and net-accumulation magnitudes of Section~\ref{subsec:regime_assignment} separate in the data: the round trips and the accumulating component. Equation~\eqref{eq:general_dgp} does not say whether the round trips belong to the signal or to the noise, and the directionality ratio has no counterpart among the model's primitives. A model with two timescales would separate the accumulating component from the round trips by horizon and give the directionality ratio as the size of the persistent component relative to the gross flow. Such a model would derive the assignment from the persistence of flow.

%%%%%%%%%%%%%%%%%%%%%%%%%%%%%%%%%%%%%%%%%%%%%%%%%%%%%%%%%%%%%%%%%%%%%%%%%%%%%%%
\section{Conclusions}
\label{sec:conclusions}
%%%%%%%%%%%%%%%%%%%%%%%%%%%%%%%%%%%%%%%%%%%%%%%%%%%%%%%%%%%%%%%%%%%%%%%%%%%%%%%

Studies normalize order flow by firm size or by trading activity according to the question that they ask, but no criterion has tied the choice to the investor class whose flow is measured. The matched filter supplies that criterion by dividing a class's flow by the carrier of the class's signal, the firm characteristic in proportion to which the flow carries the information that predicts returns, so classes whose signals have different carriers have matched filters with different denominators. In a Korean panel, capacity-bound institutional flow predicts next-day returns better under market-capitalization normalization and foreign flow better under traded-value normalization on the stock-days that both classes trade. Each class's filter is assigned from its documented scaling role and corroborated by a return-free directionality ratio. The contrast between the two classes and the advantage for capacity-bound flow, though not the advantage for foreign flow, remain significant on the disjoint pre-period and in the top \sizLargeQuintiles{} size quintiles. In Taiwan, where investment trusts tie position size to market capitalization as pension funds and insurers do, the contrast between investment-trust and foreign flow is also positive and significant, although foreign flow there predicts returns better under market-capitalization normalization. Market-capitalization normalization, which studies of informed institutional trading tend to use, therefore suits capacity-bound flow better than it suits foreign flow.


\section*{Acknowledgments}
The authors have no acknowledgments to declare.

\section*{Disclosure statement}
The authors report there are no competing interests to declare.

\section*{Declaration of generative AI use}
A large language model (Claude Opus 5, Anthropic) was used for the following tasks: (1) drafting manuscript prose and editing it for clarity and concision, (2) writing the Python code for data processing, analysis, and table generation, and (3) proposing analyses, the criteria for evaluating them, and interpretations of their results. The authors decided which proposals to adopt, and reviewed, verified, and approved all LLM-generated content. The authors take full responsibility for the study design, the interpretation of results, and the manuscript.

\section*{Funding}
This work was supported by the Korea Industrial Technology Association (KOITA) grant funded by the Korea Ministry of Science and ICT (MSIT) [grant number KOITA-2026-0002-97].

\section*{Data availability statement}
The data underlying this study comprise approximately \panelTotalObsApprox{} stock-day observations of daily order flow disaggregated by investor class, one for each common stock of KOSPI and KOSDAQ on each session in which it traded over \sampleYears{}. The pre-period estimates use the same series over \preYears{}. The flow and the total-return adjusted closes are sourced from the Korea Exchange \citep[KRX;][]{KRX2024} via the FnGuide DataGuide `all codes' export \citep{FnGuide2024}. Daily share volume, traded value, market capitalization, and the closing price at which market capitalization is formed are merged in from the \texttt{marcap} dataset \citep{marcap2024}. Access to the FnGuide DataGuide export requires a paid subscription, which also gives the book equity from which Appendix~\ref{subsec:appendix_j_factor} builds a value factor. The daily KOSPI~200 membership panel of Appendix~\ref{subsec:appendix_kospi200}, reconstructed from KRX month-end constituent snapshots and constituent-change events, is distributed with the analysis code. The Taiwan panel of Appendix~\ref{sec:taiwan} comprises \twnPanelRows{} stock-day observations of daily investor-class net trading, prices, volume, shares outstanding, and corporate-action reference prices for common stocks of the Taiwan Stock Exchange (TWSE) and the Taipei Exchange (TPEx) over \twnYears{}, retrieved from the FinMind API \citep{FinMind2024}, whose historical endpoints are openly accessible.

\section*{Code availability statement}
The analysis code that reproduces every table and reported number in
this article is available at
\url{https://github.com/hpicsk/Signal-Extraction-Order-Flow}. The repository
contains the full pipeline, from panel construction to table
generation, together with the manuscript source and a replication guide. Random
seeds are fixed in the released code, so every table and reported number
reproduces given the same data snapshot. The licensed
FnGuide data are not redistributed with the code.

%%%%%%%%%%%%%%%%%%%%%%%%%%%%%%%%%%%%%%%%%%%%%%%%%%%%%%%%%%%%%%%%%%%%%%%%%%%%%%%
% REFERENCES
%%%%%%%%%%%%%%%%%%%%%%%%%%%%%%%%%%%%%%%%%%%%%%%%%%%%%%%%%%%%%%%%%%%%%%%%%%%%%%%

\bibliographystyle{rQUF}
\bibliography{references}

\appendices
\section{Proof of Proposition~\ref{prop:matched_filter} and Corollary~\ref{cor:regimes}}
\label{sec:appendix_derivation}
\label{subsec:proof_matched_filter}

\begin{proof}
Write $\sigma_R^2 = \gamma^2\sigma_s^2 + \sigma_\epsilon^2$. Both return covariances below are positive, so the ordering of the correlations is the ordering of the squared correlations.

Because $\nu$ is independent of $(s, \omega)$ and $\epsilon$ of $s$, the matched filter's covariance with returns is $\mathrm{Cov}(S^{\text{matched}}, R) = k\gamma\sigma_s^2$, which does not depend on $\omega$. The matched filter's variance is $\mathrm{Var}(S^{\text{matched}}) = k^2\sigma_s^2 + \sigma_\nu^2 E[\omega^{-2}]$, so
\begin{equation}
\rho^2_{\text{matched}} = \frac{(k\gamma\sigma_s^2)^2}{\left(k^2\sigma_s^2 + \sigma_\nu^2 E[\omega^{-2}]\right) \sigma_R^2}.
\label{eq:snr_matched}
\end{equation}
Equation~\eqref{eq:mismatched} and the independence of $s$ from $\omega$ assumed in Remark~\ref{rem:independence} give $\mathrm{Cov}(S^{\text{mismatched}}, R) = k\gamma\sigma_s^2 E[\omega]$ and $\mathrm{Var}(S^{\text{mismatched}}) = k^2\sigma_s^2 E[\omega^2] + \sigma_\nu^2$, so
\begin{equation}
\rho^2_{\text{mismatched}} = \frac{(k\gamma\sigma_s^2)^2 (E[\omega])^2}{\left(k^2\sigma_s^2 E[\omega^2] + \sigma_\nu^2\right) \sigma_R^2}.
\label{eq:snr_mismatched}
\end{equation}
By the proof's first paragraph, the matched filter correlates better exactly when~\eqref{eq:snr_matched} exceeds~\eqref{eq:snr_mismatched}. The factor $(k\gamma\sigma_s^2)^2$ in both numerators and $\sigma_R^2$ in both denominators are strictly positive and cancel, and the two remaining denominators are strictly positive, so cross-multiplication preserves the inequality,
\begin{equation*}
k^2\sigma_s^2 E[\omega^2] + \sigma_\nu^2 > (E[\omega])^2 \left( k^2\sigma_s^2 + \sigma_\nu^2 E[\omega^{-2}] \right).
\end{equation*}
Collecting terms with the variance identity $\mathrm{Var}(\omega) = E[\omega^2] - (E[\omega])^2$ gives
\begin{equation*}
k^2\sigma_s^2 \, \mathrm{Var}(\omega) > \sigma_\nu^2 \left( (E[\omega])^2 E[\omega^{-2}] - 1 \right).
\end{equation*}
Dividing by $\sigma_\nu^2 \, \mathrm{Var}(\omega)$, which is strictly positive under the non-degeneracy assumption, gives equation~\eqref{eq:threshold}. For positive $x$ and $m$,
\begin{equation*}
\frac{1}{x^2} = \frac{1}{m^2} - \frac{2(x - m)}{m^3} + \frac{(x - m)^2 (2x + m)}{x^2 m^3}.
\end{equation*}
The ratio $\omega$ is positive because $G$ and $H$ are, so the identity holds at $x = \omega$ and $m = E[\omega]$. The identity's middle term then has expectation zero, so $E[\omega^{-2}] - (E[\omega])^{-2}$ is the expectation of the identity's last term. That expectation is positive whenever $\mathrm{Var}(\omega) > 0$, because the last term is nonnegative and is positive wherever $\omega \neq E[\omega]$. Hence $(E[\omega])^2 E[\omega^{-2}] > 1$, so $\Phi(\omega) > 0$ and the mismatched filter correlates at least as well as the matched one at every SNR in the nonempty interval $(0, \Phi(\omega)]$. Corollary~\ref{cor:regimes} follows by substituting $(G,H) = (M,V)$ and $(G,H) = (V,M)$, which set $\omega = \tau^{-1}$ and $\omega = \tau$ respectively.
\end{proof}

\section{The Korean panel: sources, universe, and sample restrictions}
\label{sec:appendix_data}

\subsection{Sources and universe}

The universe is the common stocks of KOSPI and KOSDAQ, which excludes preferred lines, SPACs, REITs, other listed investment companies, ETFs and the listings of KONEX (the exchange's market for start-up firms). The \texttt{marcap} dataset's daily snapshots give each issue's market and the sessions on which the issue is listed. The DataGuide export gives no price for a preferred line, REIT, other listed investment company or ETF, so the universe admits an issue on the sessions for which the export gives a price. SPACs, for which the export also gives prices, are excluded by their listed names. No minimum market capitalization is imposed. Turnover, the factor by which the two normalizations differ, varies widely across the panel's stock-days (mean \descTurnMean\%, standard deviation \descTurnStd\%), so across the panel the two normalized flows are far from proportional to each other.

\subsection{Investor classes}

The KRX discloses daily net trading by eight sub-types nested under the Domestic Institutional (gigwan-gye) aggregate: national pension and other public funds (yeon-gigeum-deung), insurance (boheom), investment trust (tusin), private equity funds (samo-peondeu), proprietary trading (geumyung-tuja), other financial (gita-geumyung), banks (eunhaeng), and government (gukga). Appendix~\ref{sec:appendix_subinvestor_decomposition} estimates each sub-type with nonzero flow on its own stock-days. The registered-foreigner (deungnok-oegugin) definition of foreign flow excludes the second of the two foreign categories that the exchange discloses, an unregistered other-foreigner (gita-oegugin) residual that accounts for \foreignOtherGrossPct\% of the value that foreign investors traded over the sample. Adding that residual back changes the sign of daily net foreign flow on \foreignOtherFlipPct\% of stock-days.

\subsection{Sample restrictions}
\label{subsec:appendix_panel_filters}

Of the panel's stock-days, capacity-bound institutions have nonzero net buying on \sampInstPanelPct\% and foreign investors on \sampForeignPanelPct\%. Capacity-bound institutions do not trade the stock at all on \capZeroAbsentPct\% of their stock-days with $D = 0$, and foreign investors on \foreignZeroAbsentPct\% of theirs. $D$ alone does not separate a day without trade from a day of balanced trading, so including the zero stock-days would give a day without trade the same normalized flow, zero, as a day of balanced trading.

Each daily cross-section must contain at least \minStocksPerDay{} stocks with non-missing regressors before the day's coefficient enters the time-series average, because a slope fitted to a few stocks adds mostly noise to that average. This minimum applies to every Fama-MacBeth estimate in the paper. It removes no day of the full cross-sections, whereas it removes days inside the size quintiles of Appendix~\ref{sec:appendix_liquidity_stratified}, whose day counts $T$ show how many remain.

%%%%%%%%%%%%%%%%%%%%%%%%%%%%%%%%%%%%%%%%%%%%%%%%%%%%%%%%%%%%%%%%%%%%%%%%%%%%%%%

\section{The directionality ratio and index membership}
\label{sec:appendix_scaling}

\subsection{Trading frequency and the directionality ratio}
\label{subsec:appendix_directionality_frequency}

Table~\ref{tab:directionality_frequency} reports each class's directionality ratio against the sign-flip null of Section~\ref{subsec:regime_assignment} over \dirNullDraws{} draws. Every Korean cell separates from its null at $p \leq \dirNullP$, the smallest $p$-value that the draw count allows. A closed-form null would need no simulation. Under Gaussian independent flow, however, the ratio of a directionless class is approximately the reciprocal square root of the day count, which gives the capacity-bound class \dirGaussianCap{}, whereas its sign-flip null is \dirNullCap. The difference comes from the heavy right tail of daily flow magnitude, which concentrates the numerator on a few days and raises the ratio that a directionless class attains. The Gaussian null would therefore overstate the excess directionality of both Korean classes, whereas the sign-flip null keeps the measured tail.

\begin{table}
\begin{center}
\begin{minipage}{\textwidth}
\tbl{The directionality ratio against a sign-flip null}
{\begin{tabular*}{\hsize}{\x lccccr@{}}\toprule
Investor class & Median days & Observed & Null & Excess & $N$ \\\colrule
\multicolumn{6}{@{}l@{}}{\textbf{Panel A: Korean analysis window (\sampleYears)}} \\\colrule
Capacity-bound institutions & 894 & 0.082 & 0.053 & 1.54 & 1{,}146 \\
Foreign & 1{,}471 & 0.030 & 0.044 & 0.69 & 2{,}549 \\
\colrule
\multicolumn{6}{@{}l@{}}{\textbf{Panel B: Korean pre-period (2015--2019)}} \\\colrule
Capacity-bound institutions & 764 & 0.104 & 0.052 & 2.01 & 934 \\
Foreign & 1{,}218 & 0.039 & 0.045 & 0.88 & 2{,}095 \\
\colrule
\multicolumn{6}{@{}l@{}}{\textbf{Panel C: Taiwan (2020--2025)}} \\\colrule
Investment trust (SITC) & 731 & 0.059 & 0.062 & 0.94$^{\dagger}$ & 445 \\
Foreign & 1{,}413 & 0.025 & 0.045 & 0.56 & 1{,}815 \\\botrule

\end{tabular*}}
\tabnote{Notes: `Observed' is the median across stocks of the directionality ratio $|\overline{D_t}| / \overline{|D_t|}$, the statistic of Table~\ref{tab:directionality_kr}. `Null' is its mean over \dirNullDraws{} sign-flip draws, each replacing every day's $D_t$ by $\epsilon_t |D_t|$ with $\epsilon_t = \pm 1$ independent and equally likely. `Excess' is the observed ratio divided by the null. `Median days' is the median across stocks of the class's nonzero-flow day count. Each row restricts to stocks with at least \persistMinDays{} nonzero-flow days, which for the Korean rows is the sample of Table~\ref{tab:directionality_kr}. An unmarked cell has the smallest two-sided randomization $p$-value that the draw count allows, $2/(\dirNullDraws{}+1)$, because no draw of the null reached the observed ratio, whereas a dagger marks a cell whose $p$-value exceeds that smallest value. The ratios are computed from flows in KRW for Korea and in TWD at the traded close for Taiwan.}
\label{tab:directionality_frequency}
\end{minipage}
\end{center}
\end{table}

\subsection{Foreign flow across KOSPI~200 membership}
\label{subsec:appendix_kospi200}

KOSPI~200 daily membership is reconstructed from KRX month-end snapshots and constituent-change events, with synthetic events added for unrecorded delistings and for transfers of a listing between KOSPI and KOSDAQ. Table~\ref{tab:foreign_kospi200_split} shows the predictive power of foreign $S^{\mathrm{TV}}$ outside the index, whereas passive index tracking, a channel specific to foreign flow, would concentrate that predictive power inside the index, because foreign passive funds track the MSCI Korea Standard index, whose constituents overlap closely with those of the KOSPI~200. The KOSPI~200 is a broader passively held set than the MSCI Korea Standard index, so splitting by KOSPI~200 membership dilutes a passive-tracking effect inside the in-index subset and, if anything, understates it. The foreign-flow panel of \retForeignKospiAllObs{} stock-days splits into the KOSPI~200 subset (\retForeignKospiInObs{} stock-days, \retForeignKospiInDailyAvg{} stocks per day on average) and the non-KOSPI~200 subset (\retForeignKospiOutObs{} stock-days, \retForeignKospiOutDailyAvg{} stocks per day on average). The return regressions of Table~\ref{tab:foreign_kospi200_split} use the stock-days of each that have a next-day return.

\begin{table}
\begin{center}
\begin{minipage}{\textwidth}
\tbl{Foreign flow by KOSPI~200 membership: return prediction}
{\begin{tabular*}{\hsize}{\x lcccc@{}}\toprule
Subset & $t_{\mathrm{TV}}^{\text{uni}}$ & $t_{\mathrm{TV}}^{\text{HR}}$ & $t_{\mathrm{MC}}^{\text{HR}}$ & Avg.\ $N$ stocks/day \\\colrule
Full foreign panel & 18.05*** & 18.66*** & $-1.92$* & 2,239 \\
KOSPI 200 subset & $-6.39$*** & $-4.93$*** & $-0.14$ & 200 \\
Non-KOSPI 200 subset & 20.38*** & 21.16*** & $-2.44$** & 2,039 \\
\quad largest size quintile & 4.87*** & 3.41*** & 0.78 & 253 \\\botrule

\end{tabular*}}
\tabnote{Notes: The foreign-investor flow panel partitioned by KOSPI~200 membership, with the last row restricting the non-member subset to the largest within-day market-capitalization quintile of the full cross-section. The $t$-statistics are from Fama-MacBeth regressions of next-day returns on $S^{\mathrm{TV}}$ (univariate column) and on $\{S^{\mathrm{MC}}, S^{\mathrm{TV}}\}$ jointly (horse-race columns), estimated on the same partition. ***, **, * denote significance at the 1\%, 5\%, and 10\% levels.}
\label{tab:foreign_kospi200_split}
\end{minipage}
\end{center}
\end{table}

%%%%%%%%%%%%%%%%%%%%%%%%%%%%%%%%%%%%%%%%%%%%%%%%%%%%%%%%%%%%%%%%%%%%%%%%%%%%%%%

\section{Extended robustness battery}
\label{sec:appendix_robustness_battery}
\label{sec:supp_robustness_battery}

Table~\ref{tab:robustness_summary} reports, for each of nine robustness checks, a statistic of the capacity-bound matched filter and, where the check applies to foreign flow, the same statistic of the foreign matched filter. The statistic is the filter's horse-race $t$-statistic in every check except the long-short spread and the factor-neutral alpha, which the table reports in basis points a day. The foreign matched filter's statistic is significant at the 1\% level in all seven checks that the table reports for foreign flow. The other two test $S^{\mathrm{MC}}$ alone: whether small stocks produce its result, and whether the closing price in its denominator does.

\begin{table}
\begin{center}
\begin{minipage}{\textwidth}
\tbl{Extended robustness: summary statistics for capacity-bound institutional $S^{\mathrm{MC}}$ and foreign $S^{\mathrm{TV}}$}
{\begin{tabular}{@{}lp{0.3\textwidth}p{0.3\textwidth}l@{}}\toprule
Test & Capacity-bound inst., $S^{\mathrm{MC}}$ & Foreign, $S^{\mathrm{TV}}$ & Reference \\\colrule
Pre-period hold-out (\preYears{}) & $t_{\mathrm{MC}}^{\mathrm{HR}} = \oosPPMChr$ ($t_{\mathrm{TV}}^{\mathrm{HR}} = \oosPPTVhr$) & $t_{\mathrm{TV}}^{\mathrm{HR}} = \ffPPTVhrt$ ($t_{\mathrm{MC}}^{\mathrm{HR}}$ also positive at \ffPPMChrt) & §\ref{subsec:appendix_j_oos} \\
Yearly stability & $t_{\mathrm{MC}}^{\mathrm{HR}} \geq \ryMinHRMCt$ in every year of \ryYearsCovered & $t_{\mathrm{TV}}^{\mathrm{HR}}$ from \fryMinHRTVt{} (\fryMinHRTVYear) to \fryMaxHRTVt{} (\fryMaxHRTVYear) over \fryYearsCovered & §\ref{subsec:appendix_j_oos} \\
Long-short spread, top minus bottom quintile of the filter & \lsMCbps{} bps/day ($t = \lsMCt$) & \fLsTVbps{} bps/day ($t = \fLsTVt$) & §\ref{subsec:appendix_j_longshort} \\
Winsorization (\winsorLowerPct\%/\winsorUpperPct\%) & $t_{\mathrm{MC}}^{\mathrm{HR}}$: \rpHRMCtone{} $\to$ \winMCt & $t_{\mathrm{TV}}^{\mathrm{HR}}$: \ffBaseTVhrt{} $\to$ \ffWinTVhrt{} ($t_{\mathrm{MC}}^{\mathrm{HR}}$: $\ffBaseMChrt$ $\to$ \ffWinMChrt) & §\ref{subsec:appendix_j_outlier} \\
Within-day rank transform & $t_{\mathrm{MC}}^{\mathrm{HR}}$: \rpHRMCtone{} $\to$ \rankCapHRMCt & $t_{\mathrm{TV}}^{\mathrm{HR}}$: \ffBaseTVhrt{} $\to$ \rankForHRTVt{} ($t_{\mathrm{MC}}^{\mathrm{HR}}$: $\ffBaseMChrt$ $\to$ \rankForHRMCt) & §\ref{subsec:appendix_j_rank} \\
Autocorrelation-robust (\nwLagsTen{} lags) & $t_{\mathrm{MC}}^{\mathrm{HR}} = \nwTenMCt$ & $t_{\mathrm{TV}}^{\mathrm{HR}} = \fNwTenTVt$ & §\ref{subsec:appendix_j_outlier} \\
Factor-neutral $\alpha$ (market, size, value) & $\alpha = \fnAlpha$ bps, $t = \fnAlphaT$ & $\alpha = \fFnAlpha$ bps, $t = \fFnAlphaT$ & §\ref{subsec:appendix_j_factor} \\
Market cap threshold & $t_{\mathrm{MC}}^{\mathrm{HR}} = \threshHeldMCt$ dropping the smallest \threshHeldRemovedPct\%, \threshTopMCt{} keeping the largest \threshTopTickers{} stocks & --- & §\ref{subsec:appendix_j_marketcap} \\
Formation price (VWAP cap) & $t_{\mathrm{MC}}^{\mathrm{HR}}$: \vwapCapMatchedClose{} $\to$ \vwapCapMatchedVwap & --- & §\ref{subsec:vwap_formation} \\\botrule
\end{tabular}}
\tabnote{Notes: Each row is one robustness check. Each class column reports the summary statistic for that class's return-prediction specification, estimated on the class's nonzero-flow stock-days. $t_{\mathrm{MC}}^{\mathrm{HR}}$ and $t_{\mathrm{TV}}^{\mathrm{HR}}$ are the horse-race $t$-statistics of $S^{\mathrm{MC}}$ and $S^{\mathrm{TV}}$. A dash marks a check of $S^{\mathrm{MC}}$ alone, which the table reports for capacity-bound flow only (Appendices~\ref{subsec:appendix_j_marketcap} and~\ref{subsec:vwap_formation}). Appendix~\ref{subsec:vwap_formation} reports the foreign horse race under the formation-price check as well. The autocorrelation-robust row uses Newey-West standard errors, which allow the daily slopes to be autocorrelated up to the stated number of lags. The Reference column names the subsection of Appendix~\ref{sec:appendix_robustness_battery} that holds each check's design. Of the foreign column's estimates, only the rank-transform ones also appear in a subsection (Appendix~\ref{subsec:appendix_j_rank}). Appendix~\ref{subsec:appendix_foreign_robustness} discusses the signs of the foreign pre-period coefficients.}
\label{tab:robustness_summary}
\end{minipage}
\end{center}
\end{table}

\subsection{The pre-period and yearly stability}
\label{subsec:appendix_j_oos}

The pre-period, \preYears{}, holds \sampPPDays{} daily return-prediction cross-sections that do not overlap the \sampleYears{} sample. It is a hold-out sample because the regime is assigned from each class's mandate or execution convention, and no pre-period return enters any statistic behind the primary result. The capacity-bound horse-race coefficient on the matched filter is positive at the 5\% level in each of the \ryNYears{} calendar years of the sample, from $t_{\mathrm{MC}}^{\mathrm{HR}} = \ryMinHRMCt$ in \ryMinHRMCYear{} to \ryMaxHRMCt{} in \ryMaxHRMCYear. The mismatched coefficient is negative at the 5\% level in \rySignifNegYears{} and not distinguishable from zero in \ryInsignifYears. The pre-period stacked contrast is estimated on \revPreCommonStockDays{} stock-days of common support. On the pre-period union stacking, the contrast is \revPreUnionRev{} basis points ($t = \revPreUnionT$) and the advantage for foreign flow $\revPreUnionForEdge$ ($t = \revPreUnionForEdgeT$), so that advantage is significant at the 5\% level on neither stacking.

\subsection{Long-short portfolio performance}
\label{subsec:appendix_j_longshort}

Each day we sort stocks into quintiles by a normalized flow of the capacity-bound class and hold the top quintile long and the bottom quintile short, equal-weighted and rebalanced daily. The $S^{\mathrm{MC}}$ portfolio earns \lsMCbps{} basis points a day ($t = \lsMCt$) against \lsTVbps{} basis points for the $S^{\mathrm{TV}}$ portfolio.

\subsection{Factor neutrality}
\label{subsec:appendix_j_factor}

The factors are daily analogues of the market (MKT), size (SMB), and value (HML) factors built from the same Korean universe. MKT is the value-weighted market return. SMB is the median-split size return difference, re-formed on each day's median capitalization. HML is the median-split book-to-market difference, re-formed annually at end-June following \citet{FAMA19933} on FnGuide IFRS consolidated book equity. Both legs of SMB and HML are equal-weighted. Regressed on the three factors with Newey-West standard errors \citep{NeweyWest1987}, the capacity-bound $S^{\mathrm{MC}}$ long-short portfolio earns a daily alpha of \fnAlpha{} basis points ($t = \fnAlphaT$) over \sampFactorDays{} days. Its SMB loading is negative ($\beta_{\mathrm{SMB}} = \fnSMB$, $t = \fnSMBt$), a tilt toward large stocks rather than a size premium.

\subsection{Market capitalization threshold sensitivity}
\label{subsec:appendix_j_marketcap}

The capacity-bound Fama-MacBeth analysis is re-run under six market-capitalization floors, \threshGridKRW{} billion KRW, which discard between nothing and \threshTopRemovedPct\% of stock-days. The five positive floors are rounded quantiles of the panel's capitalization distribution, and the floor of 0 is the unfiltered baseline. Both horse-race statistics clear the 1\% level, the mismatched filter's with its negative sign, up to the floor that discards \threshHeldRemovedPct\% of stock-days ($t_{\mathrm{MC}}^{\mathrm{HR}} = \threshHeldMCt$). The matched horse-race statistic attenuates monotonically across the grid and clears the 1\% level up to the floor discarding \threshMCLastPct\% of stock-days ($t_{\mathrm{MC}}^{\mathrm{HR}} = \threshMCLastT$). At the highest floor, which leaves \threshTopTickers{} stocks, neither statistic reaches the 5\% level ($t_{\mathrm{MC}}^{\mathrm{HR}} = \threshTopMCt$, $t_{\mathrm{TV}}^{\mathrm{HR}} = \threshTopTVt$). There the matched coefficient is \threshTopCoefRatio{} times its full-panel value, whereas its standard error is \threshTopSeRatio{} times as large, and the mismatched coefficient crosses zero, so the attenuation is not a precision effect alone.

Unlike the horse-race statistics, the long-short spread does not fall as the floor rises. It stays between \threshLsBpsMCMin{} and \threshLsBpsMCMax{} basis points a day on the matched filter and between \threshLsBpsTVMin{} and \threshLsBpsTVMax{} on the mismatched one. Both spreads are significant at the 1\% level at every floor. The matched filter's rank information coefficient, the mean within-day Spearman correlation with the next day's return, is indistinguishable from zero at every floor, the largest $|t|$ over the grid being \threshIcMCmaxT. Its predictive power therefore comes from the size of next-day returns rather than from their within-day ordering.

\subsection{Outlier and alternative standard errors}
\label{subsec:appendix_j_outlier}

For capacity-bound flow, winsorizing $S^{\mathrm{MC}}$, $S^{\mathrm{TV}}$, and $R_{t+1}$ within each day at the \winsorLowerPct{}st and \winsorUpperPct{}th percentiles raises the magnitude of both horse-race $t$-statistics, from $\rpHRMCtone$ to $\winMCt$ for $S^{\mathrm{MC}}$ and from $\rpHRTVtone$ to $\winTVt$ for $S^{\mathrm{TV}}$, so outliers do not produce the results. Newey-West standard errors \citep{NeweyWest1987} with \nwLagsFive{} and \nwLagsTen{} lags leave the horse-race $t$-statistics nearly unchanged, with $t = \nwFiveMCt$ and $t = \nwTenMCt$ for $S^{\mathrm{MC}}$ and $t = \nwFiveTVt$ and $t = \nwTenTVt$ for $S^{\mathrm{TV}}$.

\subsection{Within-day rank transform}
\label{subsec:appendix_j_rank}

The two normalizations are not equally exposed to a few extreme stock-days, because $|S^{\mathrm{TV}}| \leq 1$ holds by construction, whereas the bound on $|S^{\mathrm{MC}}|$ is the stock's turnover $\tau$ (Section~\ref{subsec:threshold_estimate}), which differs across stocks. We replace each normalized flow by its within-day van der Waerden score, the standard normal quantile of $r_i / (N_t+1)$, where $r_i$ is stock $i$'s rank among the day's $N_t$ stocks, and re-estimate the per-class regressions of Table~\ref{tab:investor_heterogeneity} and the stacked regressions of equation~\eqref{eq:reversal_interaction}. Within a day, the two filters then receive the same set of scores in a different order, so no difference between their slopes can come from one having the wider tail. The scores are standardized within the day as the levels are, so coefficients stay in basis points per cross-sectional standard deviation of the normalized flow. The return is left in levels, because the transform addresses the difference between the two filters' tails, and the return is the same in both regressions of a comparison.

Every univariate $t$-statistic rises in magnitude. Each assigned class keeps the larger $t$-statistic on its matched filter (Table~\ref{tab:rank_transform_class}), $t = \rankCapMCt$ on $S^{\mathrm{MC}}$ against $\rankCapTVt$ on $S^{\mathrm{TV}}$ for capacity-bound institutions and $\rankForTVt$ on $S^{\mathrm{TV}}$ against $\rankForMCt$ on $S^{\mathrm{MC}}$ for foreign investors. In the horse race, the capacity-bound mismatched filter holds its negative sign and strengthens to $\rankCapHRTVt$, whereas the foreign mismatched filter loses its negative sign and reaches only $t = \rankForHRMCt$ against $\rankForHRTVt$ on the matched filter. The contrast remains significant on both stackings, at less than half of its magnitude in levels: $\rankCommonRev$ basis points ($t = \rankCommonT$) on common support and $\rankUnionRev$ ($t = \rankUnionT$) on the union (Table~\ref{tab:rank_transform_reversal}). The advantage for capacity-bound flow remains significant, at $\rankCommonCapEdge$ basis points ($t = \rankCommonCapEdgeT$) on common support and $\rankUnionCapEdge$ ($t = \rankUnionCapEdgeT$) on the union. The advantage for foreign flow is not significant on common support, at $\rankCommonForEdge$ basis points ($t = \rankCommonForEdgeT$), whereas it is significant on the union, at $\rankUnionForEdge$ ($t = \rankUnionForEdgeT$). Each advantage on the union is the difference between the class's two univariate slopes on its own nonzero-flow stock-days, so each assigned class also keeps the larger slope on its matched filter.

The transform shrinks the advantage mechanically. The two filters are related by $S^{\mathrm{TV}} = S^{\mathrm{MC}}\tau^{-1}$ with $\tau > 0$, so a day of constant turnover would give them identical ranks and an advantage of exactly zero. On common support, the mean within-day correlation between the two standardized filters rises under the transform from \rankCorrCapZ{} to \rankCorrCapRank{} for capacity-bound institutions and from \rankCorrForZ{} to \rankCorrForRank{} for foreign investors.

\begin{table}
\begin{center}
\begin{minipage}{\textwidth}
\tbl{Return prediction by investor class: levels versus within-day ranks (Korea)}
{\begin{tabular*}{\hsize}{\x lcccccc@{}}\toprule
& \multicolumn{2}{c}{Capacity-bound inst.} & \multicolumn{2}{c}{Foreign} & \multicolumn{2}{c}{Individual} \\
\cmidrule(lr){2-3}\cmidrule(lr){4-5}\cmidrule(lr){6-7}
Specification & Level & Rank & Level & Rank & Level & Rank \\\colrule
$S^{\mathrm{MC}}$ only & 7.93 & 9.45 & 4.99 & 13.68 & $-12.06$ & $-16.26$ \\
$S^{\mathrm{TV}}$ only & 2.19 & 5.77 & 18.05 & 18.90 & $-19.16$ & $-20.01$ \\
Horse race, $S^{\mathrm{MC}}$ & 8.71 & 10.86 & $-1.92$ & 2.03 & $-8.27$ & $-7.09$ \\
Horse race, $S^{\mathrm{TV}}$ & $-5.93$ & $-9.11$ & 18.66 & 11.18 & $-12.93$ & $-6.54$ \\\botrule

\end{tabular*}}
\tabnote{Notes: Fama-MacBeth $t$-statistics on next-day returns, estimated separately for each investor class and restricted to stock-days with nonzero order flow for that class. `Level' is the specification of Table~\ref{tab:investor_heterogeneity}, which standardizes each normalized flow within the day. `Rank' replaces each normalized flow by its within-day van der Waerden normal score, the standard normal quantile of $r_i/(N_t+1)$, before the same standardization, so the two filters have identical marginal distributions within every day. The return is unchanged.}
\label{tab:rank_transform_class}
\end{minipage}
\end{center}
\end{table}

\begin{table}
\begin{center}
\begin{minipage}{\textwidth}
\tbl{The contrast under a within-day rank transform (Korea)}
{\begin{tabular*}{\hsize}{\x lcccc@{}}\toprule
& \multicolumn{2}{c}{Common support} & \multicolumn{2}{c}{Union} \\
\cmidrule(lr){2-3}\cmidrule(lr){4-5}
& Level & Rank & Level & Rank \\\colrule
Advantage of $S^{\mathrm{MC}}$ over $S^{\mathrm{TV}}$, capacity-bound & 3.63*** & 2.44*** & 3.68*** & 2.45*** \\
& (8.15) & (10.16) & (8.25) & (10.17) \\
Advantage of $S^{\mathrm{MC}}$ over $S^{\mathrm{TV}}$, foreign & $-1.37$** & 0.02 & $-4.68$*** & $-1.14$*** \\
& ($-2.27$) & (0.05) & ($-10.94$) & ($-4.14$) \\
Contrast ($b_4 - b_3$) & 5.00*** & 2.42*** & 8.36*** & 3.59*** \\
& (6.92) & (5.97) & (13.95) & (10.20) \\
\colrule
Contrast, horse race & 10.99*** & 20.80*** & 16.97*** & 26.00*** \\
& (7.60) & (8.46) & (13.92) & (11.40) \\\botrule

\end{tabular*}}
\tabnote{Notes: Fama-MacBeth estimates of the contrasts of equation~\eqref{eq:reversal_contrast}, in basis points of next-day return per standard deviation of the normalized flow, over the per-day contrast series. The first three rows are the primary specification, which is univariate. The last row is the horse race. `Level' and `Rank' are defined in the notes to Table~\ref{tab:rank_transform_class}. Common support restricts both classes to the stock-days on which both trade, and the union stacks each class on its own nonzero-flow stock-days. ***, **, * denote significance at the 1\%, 5\%, and 10\% levels.}
\label{tab:rank_transform_reversal}
\end{minipage}
\end{center}
\end{table}

\subsection{The battery on foreign flow}
\label{subsec:appendix_foreign_robustness}

On the \preYears{} pre-period, both foreign normalized flows have positive coefficients in the univariate and the horse-race specifications. The negative horse-race coefficient on foreign $S^{\mathrm{MC}}$ is therefore a property of the \sampleYears{} window rather than of foreign flow in this market throughout.
\subsection{The nonzero-flow sample restriction}
\label{subsec:appendix_zero_flow}

Without the nonzero-flow restriction of Appendix~\ref{subsec:appendix_panel_filters}, with each class's normalized flows set to zero where the class did not trade, equation~\eqref{eq:reversal_interaction} gives a contrast of \revZeroRev{} basis points ($t = \revZeroT$) on \revZeroRows{} stacked rows. Both stackings coincide on the unrestricted sample, because every row holds both classes. Relative to the union stacking, which holds the same two class panels with the restriction applied, the advantage for capacity-bound flow falls from \revUnionCapEdge{} basis points to \revZeroCapEdge{}, whereas the advantage for foreign flow stays at $\revZeroForEdge$ against $\revUnionForEdge$.

\subsection{Matched net participation}
\label{subsec:appendix_participation_matching}

The geometric mean net participation rate $|D_{i,t}|/V_{i,t}$ is \cemIntensityCap\% for capacity-bound institutions and \cemIntensityFor\% for foreign investors, a standardized class difference of \cemGapFull{} standard deviations.

We match the two classes on that rate by coarsened exact matching \citep{Iacus2012}, which bins each matching variable and re-weights the observations within each bin. Within each date, $\log \tau_{i,t}$ and $\log|S^{\mathrm{TV}}_{i,t}|$ are cut at the pooled quintiles of the two classes under the quintile grid and at their pooled deciles under the decile grid, and each (date, turnover bin, participation bin) cell is re-weighted to the smaller of its two class counts. Cells holding one class only leave the sample. The within-day standardization and the daily fit both use the cell weights. The identity $|S^{\mathrm{MC}}_{i,t}| = |S^{\mathrm{TV}}_{i,t}|\,\tau_{i,t}$ makes both magnitudes functions of the pair that the cells bin, so the matching narrows the class difference in the magnitude of each filter as well, to the extent that the bins are fine. The standardized difference on the net participation rate falls from \cemGapFull{} standard deviations to \cemGapMatched{} under the quintile grid and \cemGapDecile{} under the decile grid, against \cemBalanceBar, the standardized difference below which the matching literature treats two groups as balanced \citep{Austin2009}. The standardized difference on the magnitude of $S^{\mathrm{MC}}$, on which no cell conditions, falls with the net-participation-rate difference, from \cemMagGapFull{} to \cemMagGapMatched{} and \cemMagGapDecile{}. The quintile grid therefore reaches the threshold on the rate and exceeds it on the magnitude, whereas the decile grid is below the threshold on both.

Under both grids and on both stackings, the contrast remains significant at the 1\% level in the univariate specification and in the horse race (Table~\ref{tab:participation_matching}). The univariate contrast on the rows of common support that the quintile grid retains is \cemOverlapRev{} basis points ($t = \cemOverlapT$) at unit weights, against \cemRev{} under the cell weights and \revCommonRev{} unmatched, so the shift of the estimate comes from the re-weighting, which narrows the difference between the two classes, rather than from the dropping of single-class cells, which only shrinks the sample. The advantage for foreign flow is no longer distinguishable from zero on common support under either grid, whereas it remains distinguishable from zero on the union. Part of that advantage's loss of significance on common support is a loss of precision, because the Kish effective sample size of the re-weighted rows, the number of unit-weight rows with the same precision, is \cemEssCapPct\% of the raw capacity-bound count and \cemEssForPct\% of the raw foreign count. The point estimate shrinks toward zero as well. The design does not separate that shrinkage from the loss of precision.

\begin{table}
\begin{center}
\begin{minipage}{\textwidth}
\tbl{The contrast on matched net participation rates (Korea)}
{\begin{tabular*}{\hsize}{\x lcccccc@{}}\toprule
& \multicolumn{3}{c}{Common support} & \multicolumn{3}{c}{Union} \\
\cmidrule(lr){2-4}\cmidrule(lr){5-7}
& Unmatched & Quintile & Decile & Unmatched & Quintile & Decile \\\colrule
Advantage of $S^{\mathrm{MC}}$ over $S^{\mathrm{TV}}$, capacity-bound & 3.63*** & 3.72*** & 3.89*** & 3.68*** & 3.84*** & 3.96*** \\
& (8.15) & (7.24) & (7.17) & (8.25) & (8.21) & (8.14) \\
Advantage of $S^{\mathrm{MC}}$ over $S^{\mathrm{TV}}$, foreign & $-1.37$** & $-0.63$ & $-0.02$ & $-4.68$*** & $-2.53$*** & $-2.10$*** \\
& ($-2.27$) & ($-0.79$) & ($-0.02$) & ($-10.94$) & ($-4.56$) & ($-4.63$) \\
Contrast ($b_4 - b_3$) & 5.00*** & 4.36*** & 3.91*** & 8.36*** & 6.37*** & 6.06*** \\
& (6.92) & (4.73) & (4.25) & (13.95) & (9.01) & (9.44) \\
\colrule
Contrast, horse race & 10.99*** & 9.65*** & 9.07*** & 16.97*** & 12.54*** & 12.53*** \\
& (7.60) & (5.79) & (5.32) & (13.92) & (9.66) & (10.00) \\\botrule

\end{tabular*}}
\tabnote{Notes: Fama-MacBeth estimates of the contrasts of equation~\eqref{eq:reversal_contrast}, in basis points of next-day return per standard deviation of the normalized flow, over the per-day contrast series. The first three rows are the primary specification, which is univariate. The last row is the horse race. `Unmatched' is the specification of Table~\ref{tab:reversal_interaction}. `Quintile' and `Decile' cut $\log \tau_{i,t}$ and $\log|S^{\mathrm{TV}}_{i,t}|$ at their pooled within-date quintiles or deciles, re-weight each (date, turnover bin, participation bin) cell to the smaller of its two class counts, and re-estimate on the re-weighted rows. Common support restricts both classes to the stock-days on which both trade, and the union stacks each class on its own nonzero-flow stock-days. $t$-statistics in parentheses. ***, **, * denote significance at the 1\%, 5\%, and 10\% levels.}
\label{tab:participation_matching}
\end{minipage}
\end{center}
\end{table}

\subsection{The closing price and a VWAP market cap}
\label{subsec:vwap_formation}

Bid-ask bounce at the formation-day close $C_t$ after net buying produces a negative next-day autocorrelation, the opposite sign to the positive slopes of Table~\ref{tab:investor_heterogeneity}, so the bounce cannot produce them. Unlike the bounce, a closing price that departs from the day's other trading prices has an effect of no fixed sign on the association between $S^{\mathrm{MC}}$ and the next-day return, so the effect could favor the capacity-bound result. We re-form the market cap on the day's volume-weighted average price, $M^{\text{vwap}} = M \times (\text{VWAP}_t / C_t)$ with $\text{VWAP}_t = V_t / \mathrm{vol}_t$ and $\mathrm{vol}_t$ the day's share volume, and re-run the horse race against the same close-to-close return, because a return formed on $\text{VWAP}_t$ as well would reintroduce a shared-price term. The capacity-bound matched coefficient moves from $t = \vwapCapMatchedClose$ on the closing price to $t = \vwapCapMatchedVwap$ on VWAP, and its mismatched one from $\vwapCapMismatchClose$ to $\vwapCapMismatchVwap$. In the foreign horse race, in which $S^{\mathrm{TV}}$ is unchanged and only $S^{\mathrm{MC}}$ is re-formed, the matched coefficient moves from $t = \vwapForMatchedClose$ to $t = \vwapForMatchedVwap$, and the mismatched one from $\vwapForMismatchClose$ to $\vwapForMismatchVwap$. The reversal therefore does not come from the closing price.

\subsection{The second day after the flow}
\label{subsec:appendix_skip_day}

On the second-day return, the stacked contrast of equation~\eqref{eq:reversal_contrast} is $\revSkipCommonRev$ basis points ($t = \revSkipCommonT$) on common support and $\revSkipUnionRev$ ($t = \revSkipUnionT$) on the union stacking, against \revCommonRev{} and \revUnionRev{} on the first day. On the union stacking, the advantage for foreign flow is positive and indistinguishable from zero ($t = \revSkipUnionForEdgeT$). In the per-class regressions of Table~\ref{tab:skip_day}, $S^{\mathrm{MC}}$ has the larger horse-race $t$-statistic in both classes on the second day.

% Table: return prediction on the first and second day after the flow
\begin{table}
\begin{center}
\begin{minipage}{\textwidth}
\tbl{Return prediction on the first and second day after the flow}
{\begin{tabular*}{\hsize}{\x llcccc@{}}\toprule
& & \multicolumn{2}{c}{Univariate $t$} & \multicolumn{2}{c}{Horse-race $t$} \\
Investor class & Filter & $[t,t{+}1]$ & $[t{+}1,t{+}2]$ & $[t,t{+}1]$ & $[t{+}1,t{+}2]$ \\\colrule
Korea capacity-bound inst. & $S^{\mathrm{MC}}$ (matched) & 7.93*** & 0.49 & 8.71*** & 1.43 \\
 & $S^{\mathrm{TV}}$ (mismatched) & 2.19** & $-1.81$* & $-5.93$*** & $-2.73$*** \\
Korea Foreign & $S^{\mathrm{TV}}$ (matched) & 18.05*** & 2.43** & 18.66*** & 1.13 \\
 & $S^{\mathrm{MC}}$ (mismatched) & 4.99*** & 2.31** & $-1.92$* & 1.76* \\\botrule

\end{tabular*}}
\tabnote{Notes: Fama-MacBeth $t$-statistics for both normalizations in each class, restricted to stock-days with nonzero flow for that class. The cells run on the panel of Section~\ref{subsec:data}. The matched filter is $S^{\mathrm{MC}}$ for the capacity-bound class and $S^{\mathrm{TV}}$ for the foreign class, and the remaining normalization is labeled mismatched. The $[t,t{+}1]$ columns are the next-day return of Table~\ref{tab:investor_heterogeneity}. The $[t{+}1,t{+}2]$ columns are the one-day return over the second day after the flow. `Univariate' regresses on the stated filter alone, `Horse-race' on $S^{\mathrm{MC}}$ and $S^{\mathrm{TV}}$ jointly and reports the stated filter's coefficient. ***, **, * denote significance at the 1\%, 5\%, and 10\% levels.}
\label{tab:skip_day}
\end{minipage}
\end{center}
\end{table}

%%%%%%%%%%%%%%%%%%%%%%%%%%%%%%%%%%%%%%%%%%%%%%%%%%%%%%%%%%%%%%%%%%%%%%%%%%%%%%%

\section{The free normalization exponent}
\label{sec:supp_signal_structure}
\label{subsec:free_exponent}

The family $D/(M^{1-\theta}V^{\theta})$ of Section~\ref{subsec:threshold_estimate} equals $S^{\mathrm{MC}}\tau^{-\theta}$, with $S^{\mathrm{MC}}$ at $\theta = 0$ and $S^{\mathrm{TV}}$ at $\theta = 1$. For each class, the exponent profile is the mean over days of the cross-sectional correlation between $S(\theta)$ and the next-day return, the quantity that Proposition~\ref{prop:matched_filter} ranks, as a function of $\theta$ on a grid from $\fexGridLo$ to \fexGridHi{} in steps of \fexGridStep. The class's best exponent $\hat\theta$ maximizes the profile. Intervals resample days over \fexBootDraws{} draws.

The foreign class's best exponent exceeds the capacity-bound class's by $\fexGap$, with a day-block bootstrap interval of $\fexGapLo$ to $\fexGapHi{}$. Turnover predicts returns on its own \citep{DATAR1998203, Barinov2014}, and the factor $\tau^{-\theta}$, which both classes' normalized flows carry, brings that effect into both profiles. The comparison therefore reads the gap, which cancels the part of that effect that the two profiles share, rather than either best exponent.

In each class, the assigned endpoint reaches a larger share of the family's largest correlation than the other endpoint, in levels and on within-day ranks. For capacity-bound flow, the correlation at the market-capitalization endpoint is \fexCapAssignedPct\% of the profile's maximum, against \fexCapOtherPct\% at the traded-value endpoint. For foreign flow, the correlation at the traded-value endpoint is \fexForAssignedPct\% of the profile's maximum, against \fexForOtherPct\% at the market-capitalization endpoint. On within-day ranks, as a share of the largest correlation that the family reaches on the grid, the assigned endpoint reaches \fexRankCapAssignedPct\% against \fexRankCapOtherPct\% on capacity-bound flow and \fexRankForAssignedPct\% against \fexRankForOtherPct\% on foreign flow.

On common support, the maximum of the foreign profile falls from $\fexForRhoMax$ to $\fexCommonForRhoMax$. The gap between the two best exponents on common support is $\fexCommonGap$, with an interval of $\fexCommonGapLo$ to $\fexCommonGapHi$ that contains zero. The foreign profile on common support is therefore too flat to locate the foreign class's best exponent. The traded-value endpoint still reaches \fexCommonForAssignedPct\% of the lower maximum, against \fexCommonForOtherPct\% for the market-capitalization endpoint.

%%%%%%%%%%%%%%%%%%%%%%%%%%%%%%%%%%%%%%%%%%%%%%%%%%%%%%%%%%%%%%%%%%%%%%%%%%%%%%%

\section{Sub-investor decomposition of the KRX domestic institutional aggregate (gigwan-gye)}
\label{sec:appendix_subinvestor_decomposition}
%%%%%%%%%%%%%%%%%%%%%%%%%%%%%%%%%%%%%%%%%%%%%%%%%%%%%%%%%%%%%%%%%%%%%%%%%%%%%%%

\subsection{Aggregate alignment and the capacity-bound subset}
\label{subsec:appendix_aggregate_alignment}

At the next-day horizon, $S^{\mathrm{MC}}$ predicts better than $S^{\mathrm{TV}}$ for the KRX domestic institutional aggregate too, at univariate $t = \ihAggregateMCt$ on $S^{\mathrm{MC}}$ against $\ihAggregateTVt$ on $S^{\mathrm{TV}}$ and horse-race $t_{\mathrm{MC}}^{\mathrm{HR}} = \ihAggregateHRMCt$ against $t_{\mathrm{TV}}^{\mathrm{HR}} = \ihAggregateHRTVt$.

The capacity-bound subset pools two disclosed categories, and the KRX reports each as an aggregate with no entity-level flow, so one large fund inside the pension category could produce the class's result on its own. Insurance flow alone has an advantage of $S^{\mathrm{MC}}$ over $S^{\mathrm{TV}}$ of \doseInsAdv{} basis points at $t = \doseInsAdvT$, on \doseInsObsApprox{} stock-days, and pension flow alone \dosePenAdv{} basis points at $t = \dosePenAdvT$ on \dosePenObsApprox{}. Because the KRX assigns each institution to one of the two categories, the two estimates come from disjoint sets of institutions, so the advantage for capacity-bound flow does not depend on either category.

With Korea's investment-trust sub-type added to the class, the contrast is \revTrustRev{} basis points on common support ($t = \revTrustT$), against \revCommonRev{} without the sub-type. The advantages are \revTrustCapEdge{} basis points for the capacity-bound class ($t = \revTrustCapEdgeT$) and $\revTrustForEdge$ for foreign flow ($t = \revTrustForEdgeT$). The advantage for foreign flow, $\revCommonForEdge$ without the sub-type, moves only because the common support grows from \revCommonStockDays{} to \revTrustStockDays{} stock-days.

\subsection{The advantage across the disclosed classes}
\label{subsec:subtype_dose_response}

For each category, we estimate the univariate regressions of Section~\ref{subsec:return_prediction} on that category's own nonzero-flow panel and take the advantage of $S^{\mathrm{MC}}$ over $S^{\mathrm{TV}}$ as the Fama-MacBeth average of the per-day difference $b_{\mathrm{MC}} - b_{\mathrm{TV}}$, the quantity whose difference between the capacity-bound and foreign classes is the contrast of equation~\eqref{eq:reversal_contrast}.

All seven institutional sub-types have a positive advantage, distinguishable from zero in all but the bank sub-type ($t = \doseBankT$).

Individual flow, which the assignment leaves out (Section~\ref{subsec:return_prediction}), has an advantage of $\doseIndAdv$ basis points ($t = \doseIndAdvT$), indistinguishable from zero. Both of its normalized flows predict next-day returns negatively, at $\doseIndMC$ basis points for $S^{\mathrm{MC}}$ and $\doseIndTV$ for $S^{\mathrm{TV}}$, so the near-zero advantage records two normalizations that agree rather than two that predict nothing.

%%%%%%%%%%%%%%%%%%%%%%%%%%%%%%%%%%%%%%%%%%%%%%%%%%%%%%%%%%%%%%%%%%%%%%%%%%%%%%%

\section{The contrast across the size distribution}
\label{sec:appendix_liquidity_stratified}

We re-estimate the return prediction of Section~\ref{subsec:return_prediction} inside within-day quintiles of market capitalization (Table~\ref{tab:liquidity_stratified}). Breakpoints are formed on the full daily universe rather than within each class's nonzero-flow subset, so a quintile holds the same stocks for both classes. Each cell is the univariate Fama-MacBeth regression of $R_{t+1}$ on one day-standardized filter over the class's nonzero-flow stock-days, estimated on the days on which the bin holds at least \minStocksPerDay{} such stocks (Appendix~\ref{subsec:appendix_panel_filters}).

\begin{table}
\begin{center}
\begin{minipage}{\textwidth}
\tbl{Return prediction within market-capitalization quintiles}
{\begin{tabular*}{\hsize}{\x lcccccc@{}}\toprule
& \multicolumn{3}{c}{Capacity-bound inst.} & \multicolumn{3}{c}{Foreign} \\
\cmidrule(lr){2-4}\cmidrule(lr){5-7}
Quintile & $t(S^{\mathrm{MC}})$ & $t(S^{\mathrm{TV}})$ & $T$ & $t(S^{\mathrm{TV}})$ & $t(S^{\mathrm{MC}})$ & $T$ \\\colrule
Q1 & 0.60 & $-0.87$ & 102 & 23.17*** & 6.13*** & 1\,472 \\
Q2 & 0.53 & $-0.89$ & 675 & 16.90*** & 1.82* & 1\,472 \\
Q3 & 3.06*** & $-0.06$ & 1\,412 & 14.37*** & 0.89 & 1\,472 \\
Q4 & 6.33*** & 0.94 & 1\,472 & 11.77*** & 3.00*** & 1\,472 \\
Q5 & 5.76*** & 3.18*** & 1\,472 & 0.64 & 1.41 & 1\,472 \\\botrule

\end{tabular*}}
\tabnote{Notes: Univariate Fama-MacBeth $t$-statistics from regressing next-day returns $R_{t+1}$ on one day-standardized filter within each quintile, restricted to stock-days with nonzero flow for the investor class. Each class's matched filter is listed first, $S^{\mathrm{MC}}$ for capacity-bound institutions and $S^{\mathrm{TV}}$ for foreign flow. Quintiles are formed by market capitalization within day on the full daily cross-section. $T$ is the number of trading days on which the bin held enough nonzero-flow stocks for that class to clear the cross-sectional minimum. ***, **, * denote significance at the 1\%, 5\%, and 10\% levels.}
\label{tab:liquidity_stratified}
\end{minipage}
\end{center}
\end{table}

\subsection{Market-capitalization quintiles}
\label{subsec:appendix_size_stratified}

Capacity-bound $S^{\mathrm{MC}}$ has a larger $t$-statistic than $S^{\mathrm{TV}}$ in every quintile. In the two smallest it is indistinguishable from zero ($t = \liqCapSizeQOneT$ and $\liqCapSizeQTwoT$), because capacity-bound flow is scarce among small stocks. Only \liqCapSizeQOneDays{} of the \sampTotalDaysMinusOne{} trading days have enough nonzero-flow smallest-quintile stocks to estimate the cell. Foreign flow is measured on every day of every quintile. The $t$-statistic of foreign $S^{\mathrm{TV}}$ declines monotonically from Q1, the smallest stocks, to Q5, the largest. In the largest quintile it ($t = \liqForSizeQFiveT$) no longer exceeds that of $S^{\mathrm{MC}}$ ($t = \liqForSizeQFiveMisT$). Neither is significant at the 5\% level there, so the ordering of foreign flow's two filters is absent from the largest quintile as a whole rather than reversed.

\subsection{Three estimations inside the size distribution}
\label{subsec:appendix_size_restricted}

Common support holds \liqCommonSizeQFivePct\% of its stock-days in the largest market-capitalization quintile, against \liqForSizeQFivePct\% of the foreign class's own nonzero-flow days. Table~\ref{tab:size_restricted} reports the three estimations of Section~\ref{subsec:size_restricted} on two subsamples of common support. The bottom subsample, the bottom \sizSmallQuintiles{} quintiles, holds \sizSmallStockDays{} stock-days (\sizSmallSharePct\% of common support) at a median of \sizSmallPerDay{} per day, and the top subsample, the top \sizLargeQuintiles{} quintiles, holds \sizLargeStockDays{} stock-days at a median of \sizLargePerDay{} per day, so on the median day each subsample holds several times the minimum of \minStocksPerDay{} stocks that a daily regression needs (Appendix~\ref{subsec:appendix_panel_filters}). Each subsample is cut before the ranks, the matching cells, and the within-day standardization are formed, so each transform is formed on the cross-section on which its regression is fitted.

A control of equal sample size measures what the restriction costs in sample size alone. Each of \sizResampleDraws{} draws takes a random subsample of the common-support stock-days with the bottom subsample's count on each date, and re-estimates the advantage for foreign flow on it. The draws take stock-days rather than stacked rows, so every draw stays a common-support sample. The draws put that advantage at $\sizResampleMean$ basis points on average, with a standard deviation of \sizResampleSd{}, so the bottom subsample's $\sizSmallLvFor$ lies \sizResampleZ{} standard deviations below the draws' mean. A sample of that size alone therefore does not produce the larger advantage for foreign flow. Only \sizResampleSigPct\% of the draws make the advantage for foreign flow significant at the 5\% level in levels.

Inside the subsamples, the exponent profile of Appendix~\ref{subsec:free_exponent}, estimated in levels, does not locate the gap between the two classes' best exponents, because the gap's interval extends below zero in both, down to $\sizSmallGapLo$ in the bottom subsample and $\sizLargeGapLo$ in the top.


\begin{table}
\begin{center}
\begin{minipage}{\textwidth}
\tbl{Three estimations inside the size distribution (common support)}
{\begin{tabular*}{\hsize}{\x lccc@{}}\toprule
& \multicolumn{2}{c}{Advantage of $S^{\mathrm{MC}}$ over $S^{\mathrm{TV}}$} & \\
\cmidrule(lr){2-3}
& Capacity-bound & Foreign & Contrast \\\colrule
\multicolumn{4}{@{}l@{}}{\textbf{Panel A: bottom \sizSmallQuintiles{} market-cap quintiles}} \\
Levels & \sizSmallLvCap{} (\sizSmallLvCapT) & $\sizSmallLvFor$ ($\sizSmallLvForT$) & \sizSmallLvRev{} (\sizSmallLvRevT) \\
Within-day rank transform & \sizSmallRkCap{} (\sizSmallRkCapT) & $\sizSmallRkFor$ ($\sizSmallRkForT$) & \sizSmallRkRev{} (\sizSmallRkRevT) \\
Net-participation matching & \sizSmallCemCap{} (\sizSmallCemCapT) & $\sizSmallCemFor$ ($\sizSmallCemForT$) & \sizSmallCemRev{} (\sizSmallCemRevT) \\\colrule
\multicolumn{4}{@{}l@{}}{\textbf{Panel B: top \sizLargeQuintiles{} market-cap quintiles}} \\
Levels & \sizLargeLvCap{} (\sizLargeLvCapT) & $\sizLargeLvFor$ ($\sizLargeLvForT$) & \sizLargeLvRev{} (\sizLargeLvRevT) \\
Within-day rank transform & \sizLargeRkCap{} (\sizLargeRkCapT) & \sizLargeRkFor{} (\sizLargeRkForT) & \sizLargeRkRev{} (\sizLargeRkRevT) \\
Net-participation matching & \sizLargeCemCap{} (\sizLargeCemCapT) & $\sizLargeCemFor$ ($\sizLargeCemForT$) & \sizLargeCemRev{} (\sizLargeCemRevT) \\\botrule
\end{tabular*}}
\tabnote{Notes: Fama-MacBeth estimates of the contrasts of equation~\eqref{eq:reversal_contrast} under the univariate specification, in basis points of next-day return per standard deviation of the normalized flow, with $t$-statistics in parentheses. Every cell is estimated on common support. Panels A and B cut it at the within-day market-capitalization quintile assigned on the full daily cross-section. `Levels' is the primary specification, `Within-day rank transform' replaces each normalized flow by its van der Waerden normal score inside the subsample, and `Net-participation matching' re-weights the quintile-grid cells of Appendix~\ref{subsec:appendix_participation_matching} formed inside the subsample.}
\label{tab:size_restricted}
\end{minipage}
\end{center}
\end{table}

%%%%%%%%%%%%%%%%%%%%%%%%%%%%%%%%%%%%%%%%%%%%%%%%%%%%%%%%%%%%%%%%%%%%%%%%%%%%%%%

\section{Channels that could favor one normalization over the other}
\label{sec:appendix_nonuniform_channels}

\subsection{A lagged traded-value denominator}
\label{subsec:appendix_adv_normalisation}

The buys and sells that $D$ nets are part of the total that $V$ sums, so a class's own trading enters the denominator of $S^{\mathrm{TV}}$, an overlap between the flow and traded value. The overlap enters $S^{\mathrm{TV}}$ alone, because $D$ does not enter $M$. We rebuild the traded-value-normalized flow on the stock's average daily traded value (ADV) over the \advWindow{} panel days ending at $t-1$, and over windows of \advWindowMin{} to \advWindowMax{} days. With $\bar\tau_{i,t} = \mathrm{ADV}_{i,t-1}/M_{i,t}$, the rebuilt normalized flow equals $S^{\mathrm{MC}}_{i,t}/\bar\tau_{i,t}$, the market-cap filter divided by a turnover whose traded value the day's own trades do not enter. Requiring a full window at the longest length drops each stock's first \advWindowMax{} observations and retains \advSampleShare\% of the panel. Every window length and both denominators are estimated on that sample. Table~\ref{tab:nonuniform_channels} reports the \advWindow{}-day window.

% Table: every estimate of the reversal under the lagged traded-value denominator
\begin{table}
\begin{center}
\begin{minipage}{\textwidth}
\tbl{The contrast under a lagged traded-value denominator}
{\begin{tabular*}{\hsize}{\x lc@{}}\toprule
Estimate & Basis points ($t$) \\\colrule
Contrast, same-day denominator & \advBaseRev{} (\advBaseT) \\
Contrast, lagged denominator & \advRev{} (\advT) \\
Union stacking & \advUnionRev{} (\advUnionT) \\
Horse race & \advJointRev{} (\advJointT) \\
Horse race, cluster-by-date $t$ & \advClusterT{} \\\botrule
\end{tabular*}}
\tabnote{Notes: Each row is one estimate of the contrast of equation~\eqref{eq:reversal_contrast}, in basis points of next-day return per standard deviation of the normalized flow with $t$-statistics in parentheses. The first two rows are the univariate specification on common support, estimated on the sample with a full window, so they differ in the denominator alone. `Union stacking' stacks every stock-day either class trades, and `Horse race' enters both filters jointly. The last row refits the horse race by pooled ordinary least squares with standard errors clustered by date, absorbing the day-by-class fixed effects.}
\label{tab:nonuniform_channels}
\end{minipage}
\end{center}
\end{table}

Under the lagged denominator on common support, both classes' advantages keep the predicted signs, $\advCapEdge$ basis points for capacity-bound institutions ($t = \advCapEdgeT$) and $\advForEdge$ for foreign investors ($t = \advForEdgeT$), against \advCapEdgeBase{} and $\advForEdgeBase$ under the same-day denominator. Windows of \advWindowMin{} to \advWindowMax{} days put the contrast between \advSweepRevLo{} and \advSweepRevHi{} basis points, with $t$ of at least \advSweepTLo{}. The lagged denominator also differs from the same-day denominator in scale, because it replaces the day's traded value with its recent average rather than removing only the overlap. The gap of \advGap{} basis points between the contrasts under the two denominators therefore mixes what the overlap between $D$ and $V$ contributes to the contrast with the effect of the change of scale, so the gap measures neither alone.

\subsection{Flow persistence and the pressure it would exert}
\label{subsec:appendix_flow_continuation}

We measure the price pressure that a class exerts on day $t+1$ by its signed net participation rate $D_{i,t+1}/V_{i,t+1}$, the variable on which price impact depends in execution-cost models \citep{Almgren2001, Almgren2003}, and, as a check on that choice, by $D_{i,t+1}/M_{i,t+1}$. The pressure regressions of each class are estimated on its own nonzero-flow stock-days, the panels of Table~\ref{tab:investor_heterogeneity}, with the day-$(t+1)$ flow taken from the stock's next trading session. A class absent on day $t+1$ exerts no pressure, so its pressure variable on day $t+1$ is zero rather than missing.

Three of the four filters persist, as the channel requires: both capacity-bound filters and the foreign traded-value filter. Regressing each filter's day-$(t+1)$ value on its own day-$t$ value, both standardized within the day, gives a slope of \contCapMC{} ($t = \contCapMCT$) for the capacity-bound market-cap filter and \contCapTV{} ($t = \contCapTVT$) for its traded-value filter. Foreign flow gives $\contForMC$ ($t = \contForMCT$) for the market-cap filter and \contForTV{} ($t = \contForTVT$) for the traded-value filter, so its market-cap filter is negatively autocorrelated.

Persistence cannot produce the reversal by itself, because the pressure channel restricts the two filters jointly. Write $q_{\mathrm{MC}}$ and $q_{\mathrm{TV}}$ for the slopes of the class's day-$(t+1)$ pressure on the two filters, and $b_{\mathrm{MC}}$ and $b_{\mathrm{TV}}$ for their return slopes on the same daily cross-section. If pressure produces all of the return slope, then $b_f = \kappa_t q_f$ for each filter $f$, where $\kappa_t$ is the day's price impact per unit of pressure, taken as common to the day's stocks. Price impact does not depend on the denominator that a researcher divides by, so $\kappa_t$ is the same for both filters. Because one coefficient applies to both filters, $\kappa_t$ is over-identified, and the two filters must satisfy the restriction
\begin{equation}
b_{\mathrm{MC}}\,q_{\mathrm{TV}} - b_{\mathrm{TV}}\,q_{\mathrm{MC}} = 0,
\label{eq:proportionality}
\end{equation}
which holds for any impact coefficient, including one that changes from day to day.

The slopes of the day-$(t+1)$ net participation rate are \propCapQMC{} ($t = \propCapQMCT$) on the capacity-bound market-cap filter and \propCapQTV{} ($t = \propCapQTVT$) on its traded-value filter, and \propForQMC{} ($t = \propForQMCT$) against \propForQTV{} ($t = \propForQTVT$) for foreign flow. The traded-value slopes equal the persistence coefficients above, because the day-$(t+1)$ net participation rate is the traded-value filter's day-$(t+1)$ value. The left side of equation~\eqref{eq:proportionality} is \propCapCross{} basis points for capacity-bound flow and \propForCross{} for foreign flow, with day-block bootstrap intervals over \propBootDraws{} resamples from \propCapCrossLo{} to \propCapCrossHi{} and from \propForCrossLo{} to \propForCrossHi{} ($z = \propCapCrossZ$ and \propForCrossZ), so the restriction is rejected for both classes. With pressure measured as $D_{i,t+1}/M_{i,t+1}$, the left side is \propCapCrossCap{} and \propForCrossCap{} basis points ($z = \propCapCrossCapZ$ and \propForCrossCapZ), so the restriction is rejected again.

%%%%%%%%%%%%%%%%%%%%%%%%%%%%%%%%%%%%%%%%%%%%%%%%%%%%%%%%%%%%%%%%%%%%%%%%%%%%%%%

\section{Taiwan as a second market}
\label{sec:taiwan}

The Taiwan Stock Exchange (TWSE) and the Taipei Exchange (TPEx) disclose daily net trading by four investor classes, among them investment trusts (SITC) and foreign investors, which we retrieve through the FinMind API \citep{FinMind2024}. The panel selects common stocks by their four-digit exchange codes and holds \twnPanelRows{} stock-days of \twnPanelStocks{} such stocks over \twnYears{}, the window of the Korean panel. The exchanges report each class's net flow in shares, which we convert to TWD at the traded close. $V$ is the disclosed traded value and $M$ the traded close times shares outstanding. As in Korea, forward returns are formed on total-return adjusted closes, and sessions in which a stock did not trade are dropped. The universe is drawn from the vendor's live listing registry, which omits \twnSurvivorStocks{} stocks delisted inside the window, together \twnSurvivorSessions{} stock-days, so the Taiwan panel carries a survivorship bias of that size. A session enters the panel only if the vendor reports the stock's shares outstanding on it and on the stock's previous session, a rule that removes the sessions on which a stock traded on the Emerging Stock Board, the TPEx market on which a stock can trade before it lists.

Taiwan discloses no pension or insurance class, and SITC are regulated mutual funds, institutions that size positions to market capitalization under capacity constraints \citep{Grinold1999, LEWELLEN201162}, so they take the capacity regime and $S^{\mathrm{MC}}$. In Korea, adding the investment-trust sub-type to the capacity-bound class leaves the contrast significant at the 1\% level (Appendix~\ref{subsec:appendix_aggregate_alignment}), so a class of such funds can stand in for pension funds and insurers. Foreign flow takes the volume regime and $S^{\mathrm{TV}}$, as in Korea. The return-free corroboration of Section~\ref{subsec:regime_assignment} holds in part. Measured against each class's own sign-flip null, the directionality ratio orders SITC above foreign flow as it orders the Korean classes. In currency units, SITC does not separate from its own null (randomization $p = \dirNullTwSitcP$, the cell that a dagger marks in Panel~C of Table~\ref{tab:directionality_frequency}), unlike the Korean capacity-bound class. In share counts, the unit that the Taiwanese disclosure reports, SITC has an excess directionality of \dirNullTwSitcExcessShare{} ($p = \dirNullTwSitcPShare$) and foreign flow of \dirNullTwForExcessShare, so SITC lies significantly below its own null, on the side that marks round trips in Korea. The Taiwanese assignment therefore rests on the mandates and the ratio's ordering.

Dropping the stock-days whose forward window crosses a disclosed corporate action moves each horse-race $t$-statistic of Table~\ref{tab:twn_return_prediction} by at most \twnCorpActionDrift, so the price adjustments around corporate actions do not produce the Taiwanese results.

Table~\ref{tab:twn_contrast} estimates equation~\eqref{eq:reversal_interaction} on the two stackings of Section~\ref{subsec:reversal_test}, built from the Taiwan panel with investment trusts in the place of the capacity-bound class. On the union stacking, each class's advantage is the difference between the class's two univariate slopes in Table~\ref{tab:twn_return_prediction}, so the union row adds the $t$-statistics and no new point estimate.

\begin{table}[htbp]
\begin{center}
\begin{minipage}{\textwidth}
\tbl{Taiwan: return prediction by investor class}
{\begin{tabular*}{\hsize}{\x lccc@{}}\toprule
& MC only & TV only & Horse race \\\colrule
\multicolumn{4}{@{}l@{}}{\textbf{Panel A: Investment trust (SITC)}} \\\colrule
$S^{\mathrm{MC}}$ & 0.0017*** & --- & 0.0017*** \\
& (22.90) & & (20.57) \\
$S^{\mathrm{TV}}$ & --- & 0.0009*** & 0.0000 \\
& & (16.55) & (0.55) \\
\multicolumn{4}{@{}l@{}}{\textbf{Panel B: Foreign investor}} \\\colrule
$S^{\mathrm{MC}}$ & 0.0010*** & --- & 0.0008*** \\
& (22.09) & & (18.34) \\
$S^{\mathrm{TV}}$ & --- & 0.0008*** & 0.0004*** \\
& & (20.47) & (12.43) \\\botrule

\end{tabular*}}
\tabnote{Notes: Fama-MacBeth regressions of next-day returns on day-standardized normalized flows, estimated separately for each Taiwanese class on the class's nonzero-flow stock-days. `MC only' and `TV only' are univariate, and `Horse race' includes both normalized flows. $t$-statistics in parentheses are based on Fama-MacBeth time-series standard errors. ***, **, * denote significance at the 1\%, 5\%, and 10\% levels.}
\label{tab:twn_return_prediction}
\end{minipage}
\end{center}
\end{table}

% Table: the contrast on the Taiwan panel
\begin{table}[htbp]
\begin{center}
\begin{minipage}{\textwidth}
\tbl{Taiwan: each class's advantage and the contrast}
{\begin{tabular*}{\hsize}{\x lccc@{}}\toprule
& \multicolumn{2}{c}{Advantage of $S^{\mathrm{MC}}$ over $S^{\mathrm{TV}}$} & \\
\cmidrule(lr){2-3}
Stacking & Investment trust (SITC) & Foreign & Contrast \\\colrule
Common support & \twnCommonCapEdge{} (\twnCommonCapEdgeT) & \twnCommonForEdge{} (\twnCommonForEdgeT) & \twnCommonRev{} (\twnCommonT) \\
Union & \twnUnionCapEdge{} (\twnUnionCapEdgeT) & \twnUnionForEdge{} (\twnUnionForEdgeT) & \twnUnionRev{} (\twnUnionT) \\\botrule
\end{tabular*}}
\tabnote{Notes: Each row is one estimate of equation~\eqref{eq:reversal_interaction} on the Taiwan panel under the univariate specification, in basis points of next-day return per standard deviation of the normalized flow, with Fama-MacBeth $t$-statistics in parentheses. Investment trusts take the place of the capacity-bound class. The advantage of $S^{\mathrm{MC}}$ over $S^{\mathrm{TV}}$ is $b_1 - b_2$ for investment trusts and $(b_1 + b_3) - (b_2 + b_4)$ for foreign flow, and the contrast is their difference, defined in equation~\eqref{eq:reversal_contrast}. `Common support' restricts both classes to the stock-days on which both trade. `Union' stacks each class on its own nonzero-flow stock-days, the panels of Table~\ref{tab:twn_return_prediction}.}
\label{tab:twn_contrast}
\end{minipage}
\end{center}
\end{table}

\end{document}
